\documentclass[11pt]{report}

\usepackage{amsmath}
\usepackage[inlinechap]{settings/fncystyle}

\usepackage[tmargin=2cm,rmargin=0.5cm,lmargin=0.5cm,margin=1.5cm,bmargin=2cm,footskip=.2in]{geometry}
\usepackage{amsmath,amsfonts,amsthm,amssymb,mathtools}
\usepackage{graphicx}

%pacchetti per il font 

%\usepackage[varbb]{newpxmath}

%\usepackage{ebgaramond}
%\usepackage[ebgaramond]{newtxmath}

%\usepackage{newtxtext,newtxmath}
\usepackage[italian]{babel}


\usepackage{xfrac}
\usepackage[makeroom]{cancel}
\usepackage{mathtools}
\usepackage{enumitem}
\usepackage{hyperref,bookmark, theoremref}
\hypersetup{
	pdftitle={assignment},
	colorlinks=true, linkcolor=doc!90,
	bookmarksnumbered=true,
	bookmarksopen=true
}
\usepackage[most,many,breakable]{tcolorbox}
\usepackage{xcolor}
\usepackage{varwidth}
\usepackage{varwidth}
\usepackage{etoolbox}
%\usepackage{authblk}
\usepackage{nameref}
\usepackage{multicol,array}
\usepackage[ruled,vlined,linesnumbered]{algorithm2e}
\usepackage{comment} % enables the use of multi-line comments (\ifx \fi) 
\usepackage{import}
\usepackage{xifthen}
\usepackage{pdfpages}
\usepackage{transparent}
\usepackage{chngcntr}
\usepackage{tikz}
\usepackage{tikz-feynman}
\usepackage{tikz-3dplot}
\usepackage{titletoc}
\usepackage{pgfplots}
\usepackage{braket}
\usepackage{slashed}

\usepackage{xcolor}
\usepackage{titlesec}

\renewcommand{\chaptercolor}{doc!90}


\titleformat{\section}
  {\Large\bfseries\color{doc!70}}
  {\thesection}{1em}{#1}

\titleformat{\subsection}
  {\large\bfseries\color{doc!50}}
  {\thesubsection}{1em}{#1}



\pgfdeclarelayer{background}
\pgfdeclarelayer{foreground}
\pgfsetlayers{background,main,foreground}

\newcommand\mycommfont[1]{\footnotesize\ttfamily\textcolor{blue}{#1}}
\SetCommentSty{mycommfont}
\newcommand{\incfig}[1]{%
    \def\svgwidth{\columnwidth}
    \import{./figures/}{#1.pdf_tex}
}

\usepackage{tikzsymbols}
\tikzset{
	symbol/.style={
			draw=none,
			every to/.append style={
					edge node={node [sloped, allow upside down, auto=false]{$#1$}}}
		}
}
\tikzstyle{c} = [circle,fill=black,scale=0.5]
\tikzstyle{b} = [draw, thick, black, -]
\tikzset{
    vertex/.style={
        circle,
        draw,
        minimum size=6mm,
        inner sep=0pt
    }
}

\renewcommand\qedsymbol{$\Laughey$}

% \usepackage{xparse}
% \usepackage{pgffor}
% \usepackage{emoji}
%
% \NewDocumentCommand{\memoji}{m o}{%
%   \IfNoValueTF{#2}
%     {\ifmmode \text{\emoji{#1}} \else \emoji{#1}\fi} % Single emoji case
%     {\foreach \index in {1,...,#2}{\ifmmode \text{\emoji{#1}} \else \emoji{#1}\fi}} % Repeated emoji case
% }

%%%%%%%%%%%%%%%%%%%%%%%%%%%%%%
% SELF MADE COLORS
%%%%%%%%%%%%%%%%%%%%%%%%%%%%%%

\definecolor{doc}{HTML}{5f7d83}
\definecolor{myg}{RGB}{56, 140, 70}
\definecolor{myb}{RGB}{45, 111, 177}
\definecolor{myr}{RGB}{199, 68, 64}
\definecolor{mytheorembg}{HTML}{F2F2F9}
\definecolor{mytheoremfr}{HTML}{00007B}
\definecolor{mylemmabg}{HTML}{FFFAF8}
\definecolor{mylemmafr}{HTML}{983b0f}
\definecolor{mypropbg}{HTML}{f2fbfc}
\definecolor{mypropfr}{HTML}{191971}
\definecolor{myexamplebg}{HTML}{F2FBF8}
\definecolor{myexamplefr}{HTML}{88D6D1}
\definecolor{myexampleti}{HTML}{2A7F7F}
\definecolor{mydefinitbg}{HTML}{E5E5FF}
\definecolor{mydefinitfr}{HTML}{3F3FA3}
\definecolor{notesgreen}{RGB}{0,162,0}
\definecolor{myp}{RGB}{197, 92, 212}
\definecolor{mygr}{HTML}{2C3338}
\definecolor{myred}{RGB}{127,0,0}
\definecolor{myyellow}{RGB}{169,121,69}
\definecolor{myexercisebg}{HTML}{F2FBF8}
\definecolor{myexercisefg}{HTML}{88D6D1}

%%%%%%%%%%%%%%%%%%%%%%%%%%%%
% TCOLORBOX SETUPS
%%%%%%%%%%%%%%%%%%%%%%%%%%%%

\setlength{\parindent}{1cm}
%================================
% THEOREM BOX
%================================

\tcbuselibrary{theorems,skins,hooks}
\newtcbtheorem[number within=section]{Theorem}{Teorema}
{%
	enhanced,
	breakable,
	colback = mytheorembg,
	frame hidden,
	boxrule = 0sp,
	borderline west = {2pt}{0pt}{mytheoremfr},
	sharp corners,
	detach title,
	before upper = \tcbtitle\par\smallskip,
	coltitle = mytheoremfr,
	fonttitle = \bfseries\sffamily,
	description font = \mdseries,
	separator sign none,
	segmentation style={solid, mytheoremfr},
}
{th}

\tcbuselibrary{theorems,skins,hooks}
\newtcbtheorem[number within=chapter]{theorem}{Teorema}
{%
	enhanced,
	breakable,
	colback = mytheorembg,
	frame hidden,
	boxrule = 0sp,
	borderline west = {2pt}{0pt}{mytheoremfr},
	sharp corners,
	detach title,
	before upper = \tcbtitle\par\smallskip,
	coltitle = mytheoremfr,
	fonttitle = \bfseries\sffamily,
	description font = \mdseries,
	separator sign none,
	segmentation style={solid, mytheoremfr},
}
{th}


\tcbuselibrary{theorems,skins,hooks}
\newtcolorbox{Theoremcon}
{%
	enhanced
	,breakable
	,colback = mytheorembg
	,frame hidden
	,boxrule = 0sp
	,borderline west = {2pt}{0pt}{mytheoremfr}
	,sharp corners
	,description font = \mdseries
	,separator sign none
}

%================================
% Corollery
%================================
\tcbuselibrary{theorems,skins,hooks}
\newtcbtheorem[number within=section]{Corollary}{Corollary}
{%
	enhanced
	,breakable
	,colback = myp!10
	,frame hidden
	,boxrule = 0sp
	,borderline west = {2pt}{0pt}{myp!85!black}
	,sharp corners
	,detach title
	,before upper = \tcbtitle\par\smallskip
	,coltitle = myp!85!black
	,fonttitle = \bfseries\sffamily
	,description font = \mdseries
	,separator sign none
	,segmentation style={solid, myp!85!black}
}
{th}
\tcbuselibrary{theorems,skins,hooks}
\newtcbtheorem[number within=chapter]{corollary}{Corollary}
{%
	enhanced
	,breakable
	,colback = myp!10
	,frame hidden
	,boxrule = 0sp
	,borderline west = {2pt}{0pt}{myp!85!black}
	,sharp corners
	,detach title
	,before upper = \tcbtitle\par\smallskip
	,coltitle = myp!85!black
	,fonttitle = \bfseries\sffamily
	,description font = \mdseries
	,separator sign none
	,segmentation style={solid, myp!85!black}
}
{th}


%================================
% LEMMA
%================================

\tcbuselibrary{theorems,skins,hooks}
\newtcbtheorem[number within=section]{Lemma}{Lemma}
{%
	enhanced,
	breakable,
	colback = mylemmabg,
	frame hidden,
	boxrule = 0sp,
	borderline west = {2pt}{0pt}{mylemmafr},
	sharp corners,
	detach title,
	before upper = \tcbtitle\par\smallskip,
	coltitle = mylemmafr,
	fonttitle = \bfseries\sffamily,
	description font = \mdseries,
	separator sign none,
	segmentation style={solid, mylemmafr},
}
{th}

\tcbuselibrary{theorems,skins,hooks}
\newtcbtheorem[number within=chapter]{lemma}{lemma}
{%
	enhanced,
	breakable,
	colback = mylemmabg,
	frame hidden,
	boxrule = 0sp,
	borderline west = {2pt}{0pt}{mylemmafr},
	sharp corners,
	detach title,
	before upper = \tcbtitle\par\smallskip,
	coltitle = mylemmafr,
	fonttitle = \bfseries\sffamily,
	description font = \mdseries,
	separator sign none,
	segmentation style={solid, mylemmafr},
}
{th}

%================================
% Exercise
%================================

\tcbuselibrary{theorems,skins,hooks}
\newtcbtheorem[number within=section]{Exercise}{Esercizio}
{%
	enhanced,
	breakable,
	colback = myexercisebg,
	frame hidden,
	boxrule = 0sp,
	borderline west = {2pt}{0pt}{myexercisefg},
	sharp corners,
	detach title,
	before upper = \tcbtitle\par\smallskip,
	coltitle = myexercisefg,
	fonttitle = \bfseries\sffamily,
	description font = \mdseries,
	separator sign none,
	segmentation style={solid, myexercisefg},
}
{th}

\tcbuselibrary{theorems,skins,hooks}
\newtcbtheorem[number within=chapter]{exercise}{Esercizio}
{%
	enhanced,
	breakable,
	colback = myexercisebg,
	frame hidden,
	boxrule = 0sp,
	borderline west = {2pt}{0pt}{myexercisefg},
	sharp corners,
	detach title,
	before upper = \tcbtitle\par\smallskip,
	coltitle = myexercisefg,
	fonttitle = \bfseries\sffamily,
	description font = \mdseries,
	separator sign none,
	segmentation style={solid, myexercisefg},
}
{th}


%================================
% PROPOSITION
%================================

\tcbuselibrary{theorems,skins,hooks}
\newtcbtheorem[number within=section]{Prop}{Proposizione}
{%
	enhanced,
	breakable,
	colback = mypropbg,
	frame hidden,
	boxrule = 0sp,
	borderline west = {2pt}{0pt}{mypropfr},
	sharp corners,
	detach title,
	before upper = \tcbtitle\par\smallskip,
	coltitle = mypropfr,
	fonttitle = \bfseries\sffamily,
	description font = \mdseries,
	separator sign none,
	segmentation style={solid, mypropfr},
}
{th}

\tcbuselibrary{theorems,skins,hooks}
\newtcbtheorem[number within=chapter]{prop}{Proposizione}
{%
	enhanced,
	breakable,
	colback = mypropbg,
	frame hidden,
	boxrule = 0sp,
	borderline west = {2pt}{0pt}{mypropfr},
	sharp corners,
	detach title,
	before upper = \tcbtitle\par\smallskip,
	coltitle = mypropfr,
	fonttitle = \bfseries\sffamily,
	description font = \mdseries,
	separator sign none,
	segmentation style={solid, mypropfr},
}
{th}


%================================
% CLAIM
%================================

\tcbuselibrary{theorems,skins,hooks}
\newtcbtheorem[number within=section]{claim}{Claim}
{%
	enhanced
	,breakable
	,colback = myg!10
	,frame hidden
	,boxrule = 0sp
	,borderline west = {2pt}{0pt}{myg}
	,sharp corners
	,detach title
	,before upper = \tcbtitle\par\smallskip
	,coltitle = myg!85!black
	,fonttitle = \bfseries\sffamily
	,description font = \mdseries
	,separator sign none
	,segmentation style={solid, myg!85!black}
}
{th}



%================================
% EXAMPLE BOX
%================================

\newtcbtheorem[number within=section]{Example}{Esempio}
{%
	colback = doc!15
	,breakable
	,colframe = doc!80
	,coltitle = doc!90
	,boxrule = 1pt
	,sharp corners
	,detach title
	,before upper=\tcbtitle\par\smallskip
	,fonttitle = \bfseries
	,description font = \mdseries
	,separator sign none
	,description delimiters parenthesis
}
{ex}

\newtcbtheorem[number within=chapter]{example}{Esempio}
{%
	colback = doc!15
	,breakable
	,colframe = doc!80
	,coltitle = doc!90
	,boxrule = 1pt
	,sharp corners
	,detach title
	,before upper=\tcbtitle\par\smallskip
	,fonttitle = \bfseries
	,description font = \mdseries
	,separator sign none
	,description delimiters parenthesis
}
{ex}

%================================
% DEFINITION BOX
%================================

\newtcbtheorem[number within=section]{Definition}{Definizione}{enhanced,
	before skip=2mm,after skip=2mm, colback=red!5,colframe=red!80!black,boxrule=0.5mm,
	attach boxed title to top left={xshift=1cm,yshift*=1mm-\tcboxedtitleheight}, varwidth boxed title*=-3cm,
	boxed title style={frame code={
					\path[fill=tcbcolback]
					([yshift=-1mm,xshift=-1mm]frame.north west)
					arc[start angle=0,end angle=180,radius=1mm]
					([yshift=-1mm,xshift=1mm]frame.north east)
					arc[start angle=180,end angle=0,radius=1mm];
					\path[left color=tcbcolback!60!black,right color=tcbcolback!60!black,
						middle color=tcbcolback!80!black]
					([xshift=-2mm]frame.north west) -- ([xshift=2mm]frame.north east)
					[rounded corners=1mm]-- ([xshift=1mm,yshift=-1mm]frame.north east)
					-- (frame.south east) -- (frame.south west)
					-- ([xshift=-1mm,yshift=-1mm]frame.north west)
					[sharp corners]-- cycle;
				},interior engine=empty,
		},
	fonttitle=\bfseries,
	title={#2},#1}{def}
\newtcbtheorem[number within=chapter]{definition}{Definizione}{enhanced,
	before skip=2mm,after skip=2mm, colback=red!5,colframe=red!80!black,boxrule=0.5mm,
	attach boxed title to top left={xshift=1cm,yshift*=1mm-\tcboxedtitleheight}, varwidth boxed title*=-3cm,
	boxed title style={frame code={
					\path[fill=tcbcolback]
					([yshift=-1mm,xshift=-1mm]frame.north west)
					arc[start angle=0,end angle=180,radius=1mm]
					([yshift=-1mm,xshift=1mm]frame.north east)
					arc[start angle=180,end angle=0,radius=1mm];
					\path[left color=tcbcolback!60!black,right color=tcbcolback!60!black,
						middle color=tcbcolback!80!black]
					([xshift=-2mm]frame.north west) -- ([xshift=2mm]frame.north east)
					[rounded corners=1mm]-- ([xshift=1mm,yshift=-1mm]frame.north east)
					-- (frame.south east) -- (frame.south west)
					-- ([xshift=-1mm,yshift=-1mm]frame.north west)
					[sharp corners]-- cycle;
				},interior engine=empty,
		},
	fonttitle=\bfseries,
	title={#2},#1}{def}


%================================
% EXERCISE BOX
%================================

\newcounter{questioncounter}
\counterwithin{questioncounter}{chapter}
% \counterwithin{questioncounter}{section}

\makeatletter
\newtcbtheorem[use counter=questioncounter]{question}{Domanda}{enhanced,
	breakable,
	colback=white,
	colframe=myb!80!black,
	attach boxed title to top left={yshift*=-\tcboxedtitleheight},
	fonttitle=\bfseries,
	title={#2},
	boxed title size=title,
	boxed title style={%
			sharp corners,
			rounded corners=northwest,
			colback=tcbcolframe,
			boxrule=0pt,
		},
	underlay boxed title={%
			\path[fill=tcbcolframe] (title.south west)--(title.south east)
			to[out=0, in=180] ([xshift=5mm]title.east)--
			(title.center-|frame.east)
			[rounded corners=\kvtcb@arc] |-
			(frame.north) -| cycle;
		},
	#1
}{def}
\makeatother

%================================
% SOLUTION BOX
%================================

\makeatletter
\newtcolorbox{solution}{enhanced,
	breakable,
	colback=white,
	colframe=myg!80!black,
	attach boxed title to top left={yshift*=-\tcboxedtitleheight},
	title=Solution,
	boxed title size=title,
	boxed title style={%
			sharp corners,
			rounded corners=northwest,
			colback=tcbcolframe,
			boxrule=0pt,
		},
	underlay boxed title={%
			\path[fill=tcbcolframe] (title.south west)--(title.south east)
			to[out=0, in=180] ([xshift=5mm]title.east)--
			(title.center-|frame.east)
			[rounded corners=\kvtcb@arc] |-
			(frame.north) -| cycle;
		},
}
\makeatother

%================================
% Question BOX
%================================

\makeatletter
\newtcbtheorem{qstion}{Domanda}{enhanced,
	breakable,
	colback=white,
	colframe=mygr,
	attach boxed title to top left={yshift*=-\tcboxedtitleheight},
	fonttitle=\bfseries,
	title={#2},
	boxed title size=title,
	boxed title style={%
			sharp corners,
			rounded corners=northwest,
			colback=tcbcolframe,
			boxrule=0pt,
		},
	underlay boxed title={%
			\path[fill=tcbcolframe] (title.south west)--(title.south east)
			to[out=0, in=180] ([xshift=5mm]title.east)--
			(title.center-|frame.east)
			[rounded corners=\kvtcb@arc] |-
			(frame.north) -| cycle;
		},
	#1
}{def}
\makeatother

\newtcbtheorem[number within=chapter]{wconc}{Wrong Concept}{
	breakable,
	enhanced,
	colback=white,
	colframe=myr,
	arc=0pt,
	outer arc=0pt,
	fonttitle=\bfseries\sffamily\large,
	colbacktitle=myr,
	attach boxed title to top left={},
	boxed title style={
			enhanced,
			skin=enhancedfirst jigsaw,
			arc=3pt,
			bottom=0pt,
			interior style={fill=myr}
		},
	#1
}{def}


%================================
% NOTE BOX
%================================

\usetikzlibrary{arrows,calc,shadows.blur, backgrounds}
\tcbuselibrary{skins}
\newtcolorbox{note}[1][]{%
	enhanced jigsaw,
	colback=gray!20!white,%
	colframe=gray!80!black,
	size=small,
	boxrule=1pt,
	title=\textbf{Osservazione:-},
	halign title=flush center,
	coltitle=black,
	breakable,
	drop shadow=black!50!white,
	attach boxed title to top left={xshift=1cm,yshift=-\tcboxedtitleheight/2,yshifttext=-\tcboxedtitleheight/2},
	minipage boxed title=3cm,
	boxed title style={%
			colback=white,
			size=fbox,
			boxrule=1pt,
			boxsep=2pt,
			underlay={%
					\coordinate (dotA) at ($(interior.west) + (-0.5pt,0)$);
					\coordinate (dotB) at ($(interior.east) + (0.5pt,0)$);
					\begin{scope}
						\clip (interior.north west) rectangle ([xshift=3ex]interior.east);
						\filldraw [white, blur shadow={shadow opacity=60, shadow yshift=-.75ex}, rounded corners=2pt] (interior.north west) rectangle (interior.south east);
					\end{scope}
					\begin{scope}[gray!80!black]
						\fill (dotA) circle (2pt);
						\fill (dotB) circle (2pt);
					\end{scope}
				},
		},
	#1,
}

%%%%%%%%%%%%%%%%%%%%%%%%%%%%%%
% SELF MADE COMMANDS
%%%%%%%%%%%%%%%%%%%%%%%%%%%%%%

\newcommand{\thm}[2]{\begin{Theorem}{#1}{}#2\end{Theorem}}
\newcommand{\cor}[2]{\begin{Corollary}{#1}{}#2\end{Corollary}}
\newcommand{\mlemma}[2]{\begin{Lemma}{#1}{}#2\end{Lemma}}
\newcommand{\mer}[2]{\begin{Exercise}{#1}{}#2\end{Exercise}}
\newcommand{\mprop}[2]{\begin{Prop}{#1}{}#2\end{Prop}}
\newcommand{\clm}[3]{\begin{claim}{#1}{#2}#3\end{claim}}
\newcommand{\wc}[2]{\begin{wconc}{#1}{}\setlength{\parindent}{1cm}#2\end{wconc}}
\newcommand{\thmcon}[1]{\begin{Theoremcon}{#1}\end{Theoremcon}}
\newcommand{\ex}[2]{\begin{Example}{#1}{}#2\end{Example}}
\newcommand{\dfn}[2]{\begin{Definition}[colbacktitle=red!75!black]{#1}{}#2\end{Definition}}
\newcommand{\dfnc}[2]{\begin{definition}[colbacktitle=red!75!black]{#1}{}#2\end{definition}}
\newcommand{\qs}[2]{\begin{question}{#1}{}#2\end{question}}
\newcommand{\pf}[2]{\begin{myproof}[#1]#2\end{myproof}}
\newcommand{\nt}[1]{\begin{note}#1\end{note}}

\newcommand*\circled[1]{\tikz[baseline=(char.base)]{
		\node[shape=circle,draw,inner sep=1pt] (char) {#1};}}
\newcommand\getcurrentref[1]{%
	\ifnumequal{\value{#1}}{0}
	{??}
	{\the\value{#1}}%
}
\newcommand{\getCurrentSectionNumber}{\getcurrentref{section}}
\newenvironment{myproof}[1][\proofname]{%
	\proof[\bfseries #1: ]%
}{\endproof}

\newcommand{\mclm}[2]{\begin{myclaim}[#1]#2\end{myclaim}}
\newenvironment{myclaim}[1][\claimname]{\proof[\bfseries #1: ]}{}
\newenvironment{iclaim}[1][\claimname]{\bfseries #1\mdseries:}{}
\newcommand{\iclm}[2]{\begin{iclaim}[#1]#2\end{iclaim}}

\newcounter{mylabelcounter}

\makeatletter
\newcommand{\setword}[2]{%
	\phantomsection
	#1\def\@currentlabel{\unexpanded{#1}}\label{#2}%
}
\makeatother

% deliminators
\DeclarePairedDelimiter{\abs}{\lvert}{\rvert}
\DeclarePairedDelimiter{\norm}{\lVert}{\rVert}

\DeclarePairedDelimiter{\ceil}{\lceil}{\rceil}
\DeclarePairedDelimiter{\floor}{\lfloor}{\rfloor}
\DeclarePairedDelimiter{\round}{\lfloor}{\rceil}

\newsavebox\diffdbox
\newcommand{\slantedromand}{{\mathpalette\makesl{d}}}
\newcommand{\makesl}[2]{%
\begingroup
\sbox{\diffdbox}{$\mathsurround=0pt#1\mathrm{#2}$}%
\pdfsave
\pdfsetmatrix{1 0 0.2 1}%
\rlap{\usebox{\diffdbox}}%
\pdfrestore
\hskip\wd\diffdbox
\endgroup
}
\newcommand{\dd}[1][]{\ensuremath{\mathop{}\!\ifstrempty{#1}{%
\slantedromand\@ifnextchar^{\hspace{0.2ex}}{\hspace{0.1ex}}}%
{\slantedromand\hspace{0.2ex}^{#1}}}}
\ProvideDocumentCommand\dv{o m g}{%
  \ensuremath{%
    \IfValueTF{#3}{%
      \IfNoValueTF{#1}{%
        \frac{\dd #2}{\dd #3}%
      }{%
        \frac{\dd^{#1} #2}{\dd #3^{#1}}%
      }%
    }{%
      \IfNoValueTF{#1}{%
        \frac{\dd}{\dd #2}%
      }{%
        \frac{\dd^{#1}}{\dd #2^{#1}}%
      }%
    }%
  }%
}
\providecommand*{\pdv}[3][]{\frac{\partial^{#1}#2}{\partial#3^{#1}}}
%  - others
\DeclareMathOperator{\Lap}{\mathcal{L}}
\DeclareMathOperator{\Var}{Var} % varience
\DeclareMathOperator{\Cov}{Cov} % covarience
\DeclareMathOperator{\E}{E} % expected

% Since the amsthm package isn't loaded

% I prefer the slanted \leq
\let\oldleq\leq % save them in case they're every wanted
\let\oldgeq\geq
\renewcommand{\leq}{\leqslant}
\renewcommand{\geq}{\geqslant}

%%%%%%%%%%%%%%%%%%%%%%%%%%%%%%%%%%%%%%%%%%%
% TABLE OF CONTENTS
%%%%%%%%%%%%%%%%%%%%%%%%%%%%%%%%%%%%%%%%%%%

\contentsmargin{0cm}
\titlecontents{chapter}[3.7pc]
{\addvspace{30pt}%
	\begin{tikzpicture}[remember picture, overlay]%
		\draw[fill=doc!60,draw=doc!60] (-7,-.1) rectangle (0,.5);%
		\pgftext[left,x=-3.0cm,y=0.2cm]{\color{white}\Large\sc\bfseries capitolo\ \thecontentslabel};
	\end{tikzpicture}\color{doc!60}\large\sc\bfseries}%
{}
{}
{\;\titlerule\;\large\sc\bfseries Page \thecontentspage
	\begin{tikzpicture}[remember picture, overlay]
		\draw[fill=doc!60,draw=doc!60] (2pt,0) rectangle (4,0.1pt);
	\end{tikzpicture}}%
\titlecontents{section}[3.7pc]
{\addvspace{2pt}}
{\contentslabel[\thecontentslabel]{2pc}}
{}
{\hfill\small \thecontentspage}
[]
\titlecontents*{subsection}[3.7pc]
{\addvspace{-1pt}\small}
{}
{}
{\ --- \small\thecontentspage}
[ \textbullet\ ][]

\makeatletter
\renewcommand{\tableofcontents}{%
	\chapter*{%
	  \vspace*{-20\p@}%
	  \begin{tikzpicture}[remember picture, overlay]%
		  \pgftext[right,x=15cm,y=0.2cm]{\color{doc!60}\Huge\sc\bfseries \contentsname};%
		  \draw[fill=doc!60,draw=doc!60] (13,-.75) rectangle (20,1);%
		  \clip (13,-.75) rectangle (20,1);
		  \pgftext[right,x=15cm,y=0.2cm]{\color{white}\Huge\sc\bfseries \contentsname};%
	  \end{tikzpicture}}%
	\@starttoc{toc}}
\makeatother
\newcommand{\eps}{\epsilon}
\newcommand{\veps}{\varepsilon}
\newcommand{\Qed}{\begin{flushright}\qed\end{flushright}}

\newcommand{\parinn}{\setlength{\parindent}{1cm}}
\newcommand{\parinf}{\setlength{\parindent}{0cm}}

% \newcommand{\norm}{\|\cdot\|}
\newcommand{\inorm}{\norm_{\infty}}
\newcommand{\opensets}{\{V_{\alpha}\}_{\alpha\in I}}
\newcommand{\oset}{V_{\alpha}}
\newcommand{\opset}[1]{V_{\alpha_{#1}}}
\newcommand{\lub}{\text{lub}}
\newcommand{\del}[2]{\frac{\partial #1}{\partial #2}}
\newcommand{\Del}[3]{\frac{\partial^{#1} #2}{\partial^{#1} #3}}
\newcommand{\deld}[2]{\dfrac{\partial #1}{\partial #2}}
\newcommand{\Deld}[3]{\dfrac{\partial^{#1} #2}{\partial^{#1} #3}}
\newcommand{\der}[2]{\frac{\mathrm{d} #1}{\mathrm{d} #2}}
% \newcommand{\ddd}[3]{\frac{\mathrm{d}^{#3} #1}{\mathrm{d}^{#3} #2}}
\newcommand{\lm}{\lambda}
\newcommand{\uin}{\mathbin{\rotatebox[origin=c]{90}{$\in$}}}
\newcommand{\usubset}{\mathbin{\rotatebox[origin=c]{90}{$\subset$}}}
\newcommand{\lt}{\left}
\newcommand{\rt}{\right}
\newcommand{\bs}[1]{\boldsymbol{#1}}
\newcommand{\exs}{\exists}
\newcommand{\st}{\strut}
\newcommand{\dps}[1]{\displaystyle{#1}}
\newcommand{\id}{\text{id}}


\newcommand{\sol}{\setlength{\parindent}{0cm}\textbf{\textit{Solution:}}\setlength{\parindent}{1cm} }
\newcommand{\solve}[1]{\setlength{\parindent}{0cm}\textbf{\textit{Solution: }}\setlength{\parindent}{1cm}#1 \Qed}

% number sets
\newcommand{\RR}[1][]{\ensuremath{\ifstrempty{#1}{\mathbb{R}}{\mathbb{R}^{#1}}}}
\newcommand{\NN}[1][]{\ensuremath{\ifstrempty{#1}{\mathbb{N}}{\mathbb{N}^{#1}}}}
\newcommand{\ZZ}[1][]{\ensuremath{\ifstrempty{#1}{\mathbb{Z}}{\mathbb{Z}^{#1}}}}
\newcommand{\QQ}[1][]{\ensuremath{\ifstrempty{#1}{\mathbb{Q}}{\mathbb{Q}^{#1}}}}
\newcommand{\CC}[1][]{\ensuremath{\ifstrempty{#1}{\mathbb{C}}{\mathbb{C}^{#1}}}}
\newcommand{\PP}[1][]{\ensuremath{\ifstrempty{#1}{\mathbb{P}}{\mathbb{P}^{#1}}}}
\newcommand{\HH}[1][]{\ensuremath{\ifstrempty{#1}{\mathbb{H}}{\mathbb{H}^{#1}}}}
\newcommand{\FF}[1][]{\ensuremath{\ifstrempty{#1}{\mathbb{F}}{\mathbb{F}^{#1}}}}
% expected value
\newcommand{\EE}{\ensuremath{\mathbb{E}}}

%---------------------------------------
% BlackBoard Math Fonts :-
%---------------------------------------

%Captital Letters
\newcommand{\bbA}{\mathbb{A}}	\newcommand{\bbB}{\mathbb{B}}
\newcommand{\bbC}{\mathbb{C}}	\newcommand{\bbD}{\mathbb{D}}
\newcommand{\bbE}{\mathbb{E}}	\newcommand{\bbF}{\mathbb{F}}
\newcommand{\bbG}{\mathbb{G}}	\newcommand{\bbH}{\mathbb{H}}
\newcommand{\bbI}{\mathbb{I}}	\newcommand{\bbJ}{\mathbb{J}}
\newcommand{\bbK}{\mathbb{K}}	\newcommand{\bbL}{\mathbb{L}}
\newcommand{\bbM}{\mathbb{M}}	\newcommand{\bbN}{\mathbb{N}}
\newcommand{\bbO}{\mathbb{O}}	\newcommand{\bbP}{\mathbb{P}}
\newcommand{\bbQ}{\mathbb{Q}}	\newcommand{\bbR}{\mathbb{R}}
\newcommand{\bbS}{\mathbb{S}}	\newcommand{\bbT}{\mathbb{T}}
\newcommand{\bbU}{\mathbb{U}}	\newcommand{\bbV}{\mathbb{V}}
\newcommand{\bbW}{\mathbb{W}}	\newcommand{\bbX}{\mathbb{X}}
\newcommand{\bbY}{\mathbb{Y}}	\newcommand{\bbZ}{\mathbb{Z}}

%---------------------------------------
% MathCal Fonts :-
%---------------------------------------

%Captital Letters
\newcommand{\mcA}{\mathcal{A}}	\newcommand{\mcB}{\mathcal{B}}
\newcommand{\mcC}{\mathcal{C}}	\newcommand{\mcD}{\mathcal{D}}
\newcommand{\mcE}{\mathcal{E}}	\newcommand{\mcF}{\mathcal{F}}
\newcommand{\mcG}{\mathcal{G}}	\newcommand{\mcH}{\mathcal{H}}
\newcommand{\mcI}{\mathcal{I}}	\newcommand{\mcJ}{\mathcal{J}}
\newcommand{\mcK}{\mathcal{K}}	\newcommand{\mcL}{\mathcal{L}}
\newcommand{\mcM}{\mathcal{M}}	\newcommand{\mcN}{\mathcal{N}}
\newcommand{\mcO}{\mathcal{O}}	\newcommand{\mcP}{\mathcal{P}}
\newcommand{\mcQ}{\mathcal{Q}}	\newcommand{\mcR}{\mathcal{R}}
\newcommand{\mcS}{\mathcal{S}}	\newcommand{\mcT}{\mathcal{T}}
\newcommand{\mcU}{\mathcal{U}}	\newcommand{\mcV}{\mathcal{V}}
\newcommand{\mcW}{\mathcal{W}}	\newcommand{\mcX}{\mathcal{X}}
\newcommand{\mcY}{\mathcal{Y}}	\newcommand{\mcZ}{\mathcal{Z}}



%---------------------------------------
% Bold Math Fonts :-
%---------------------------------------

%Captital Letters
\newcommand{\bmA}{\boldsymbol{A}}	\newcommand{\bmB}{\boldsymbol{B}}
\newcommand{\bmC}{\boldsymbol{C}}	\newcommand{\bmD}{\boldsymbol{D}}
\newcommand{\bmE}{\boldsymbol{E}}	\newcommand{\bmF}{\boldsymbol{F}}
\newcommand{\bmG}{\boldsymbol{G}}	\newcommand{\bmH}{\boldsymbol{H}}
\newcommand{\bmI}{\boldsymbol{I}}	\newcommand{\bmJ}{\boldsymbol{J}}
\newcommand{\bmK}{\boldsymbol{K}}	\newcommand{\bmL}{\boldsymbol{L}}
\newcommand{\bmM}{\boldsymbol{M}}	\newcommand{\bmN}{\boldsymbol{N}}
\newcommand{\bmO}{\boldsymbol{O}}	\newcommand{\bmP}{\boldsymbol{P}}
\newcommand{\bmQ}{\boldsymbol{Q}}	\newcommand{\bmR}{\boldsymbol{R}}
\newcommand{\bmS}{\boldsymbol{S}}	\newcommand{\bmT}{\boldsymbol{T}}
\newcommand{\bmU}{\boldsymbol{U}}	\newcommand{\bmV}{\boldsymbol{V}}
\newcommand{\bmW}{\boldsymbol{W}}	\newcommand{\bmX}{\boldsymbol{X}}
\newcommand{\bmY}{\boldsymbol{Y}}	\newcommand{\bmZ}{\boldsymbol{Z}}
%Small Letters
\newcommand{\bma}{\boldsymbol{a}}	\newcommand{\bmb}{\boldsymbol{b}}
\newcommand{\bmc}{\boldsymbol{c}}	\newcommand{\bmd}{\boldsymbol{d}}
\newcommand{\bme}{\boldsymbol{e}}	\newcommand{\bmf}{\boldsymbol{f}}
\newcommand{\bmg}{\boldsymbol{g}}	\newcommand{\bmh}{\boldsymbol{h}}
\newcommand{\bmi}{\boldsymbol{i}}	\newcommand{\bmj}{\boldsymbol{j}}
\newcommand{\bmk}{\boldsymbol{k}}	\newcommand{\bml}{\boldsymbol{l}}
\newcommand{\bmm}{\boldsymbol{m}}	\newcommand{\bmn}{\boldsymbol{n}}
\newcommand{\bmo}{\boldsymbol{o}}	\newcommand{\bmp}{\boldsymbol{p}}
\newcommand{\bmq}{\boldsymbol{q}}	\newcommand{\bmr}{\boldsymbol{r}}
\newcommand{\bms}{\boldsymbol{s}}	\newcommand{\bmt}{\boldsymbol{t}}
\newcommand{\bmu}{\boldsymbol{u}}	\newcommand{\bmv}{\boldsymbol{v}}
\newcommand{\bmw}{\boldsymbol{w}}	\newcommand{\bmx}{\boldsymbol{x}}
\newcommand{\bmy}{\boldsymbol{y}}	\newcommand{\bmz}{\boldsymbol{z}}

%---------------------------------------
% Scr Math Fonts :-
%---------------------------------------

\newcommand{\sA}{{\mathscr{A}}}   \newcommand{\sB}{{\mathscr{B}}}
\newcommand{\sC}{{\mathscr{C}}}   \newcommand{\sD}{{\mathscr{D}}}
\newcommand{\sE}{{\mathscr{E}}}   \newcommand{\sF}{{\mathscr{F}}}
\newcommand{\sG}{{\mathscr{G}}}   \newcommand{\sH}{{\mathscr{H}}}
\newcommand{\sI}{{\mathscr{I}}}   \newcommand{\sJ}{{\mathscr{J}}}
\newcommand{\sK}{{\mathscr{K}}}   \newcommand{\sL}{{\mathscr{L}}}
\newcommand{\sM}{{\mathscr{M}}}   \newcommand{\sN}{{\mathscr{N}}}
\newcommand{\sO}{{\mathscr{O}}}   \newcommand{\sP}{{\mathscr{P}}}
\newcommand{\sQ}{{\mathscr{Q}}}   \newcommand{\sR}{{\mathscr{R}}}
\newcommand{\sS}{{\mathscr{S}}}   \newcommand{\sT}{{\mathscr{T}}}
\newcommand{\sU}{{\mathscr{U}}}   \newcommand{\sV}{{\mathscr{V}}}
\newcommand{\sW}{{\mathscr{W}}}   \newcommand{\sX}{{\mathscr{X}}}
\newcommand{\sY}{{\mathscr{Y}}}   \newcommand{\sZ}{{\mathscr{Z}}}


%---------------------------------------
% Math Fraktur Font
%---------------------------------------

%Captital Letters
\newcommand{\mfA}{\mathfrak{A}}	\newcommand{\mfB}{\mathfrak{B}}
\newcommand{\mfC}{\mathfrak{C}}	\newcommand{\mfD}{\mathfrak{D}}
\newcommand{\mfE}{\mathfrak{E}}	\newcommand{\mfF}{\mathfrak{F}}
\newcommand{\mfG}{\mathfrak{G}}	\newcommand{\mfH}{\mathfrak{H}}
\newcommand{\mfI}{\mathfrak{I}}	\newcommand{\mfJ}{\mathfrak{J}}
\newcommand{\mfK}{\mathfrak{K}}	\newcommand{\mfL}{\mathfrak{L}}
\newcommand{\mfM}{\mathfrak{M}}	\newcommand{\mfN}{\mathfrak{N}}
\newcommand{\mfO}{\mathfrak{O}}	\newcommand{\mfP}{\mathfrak{P}}
\newcommand{\mfQ}{\mathfrak{Q}}	\newcommand{\mfR}{\mathfrak{R}}
\newcommand{\mfS}{\mathfrak{S}}	\newcommand{\mfT}{\mathfrak{T}}
\newcommand{\mfU}{\mathfrak{U}}	\newcommand{\mfV}{\mathfrak{V}}
\newcommand{\mfW}{\mathfrak{W}}	\newcommand{\mfX}{\mathfrak{X}}
\newcommand{\mfY}{\mathfrak{Y}}	\newcommand{\mfZ}{\mathfrak{Z}}
%Small Letters
\newcommand{\mfa}{\mathfrak{a}}	\newcommand{\mfb}{\mathfrak{b}}
\newcommand{\mfc}{\mathfrak{c}}	\newcommand{\mfd}{\mathfrak{d}}
\newcommand{\mfe}{\mathfrak{e}}	\newcommand{\mff}{\mathfrak{f}}
\newcommand{\mfg}{\mathfrak{g}}	\newcommand{\mfh}{\mathfrak{h}}
\newcommand{\mfi}{\mathfrak{i}}	\newcommand{\mfj}{\mathfrak{j}}
\newcommand{\mfk}{\mathfrak{k}}	\newcommand{\mfl}{\mathfrak{l}}
\newcommand{\mfm}{\mathfrak{m}}	\newcommand{\mfn}{\mathfrak{n}}
\newcommand{\mfo}{\mathfrak{o}}	\newcommand{\mfp}{\mathfrak{p}}
\newcommand{\mfq}{\mathfrak{q}}	\newcommand{\mfr}{\mathfrak{r}}
\newcommand{\mfs}{\mathfrak{s}}	\newcommand{\mft}{\mathfrak{t}}
\newcommand{\mfu}{\mathfrak{u}}	\newcommand{\mfv}{\mathfrak{v}}
\newcommand{\mfw}{\mathfrak{w}}	\newcommand{\mfx}{\mathfrak{x}}
\newcommand{\mfy}{\mathfrak{y}}	\newcommand{\mfz}{\mathfrak{z}}


\usepackage[indent=0pt]{parskip}

\renewcommand{\chaptername}{Capitolo}
\renewcommand{\contentsname}{Indice}

\title{\Huge{\textbf{Master thesis notes}}}
\author{Giacomo Galliani}
\date{Academic year 2026-2027}

\begin{document}

\maketitle
\newpage% or \cleardoublepage
% \pdfbookmark[<level>]{<title>}{<dest>}
\pdfbookmark[section]{\contentsname}{toc}
\tableofcontents
\pagebreak

\chapter{NLO QCD computations in the subtraction scheme}

\section{The singularities of the NLO real corrections for DY processes}

Consider the production of a lepton pair, by annichilation of a quark and a antiquark. These process are known as DY process. At LO (in \(\alpha_s\)), there are two diagrams, when the exchange is mediated by a photon or a Z boson. The amplitude is known. We label \(p_1\) the momentum of the quark, and \(p_2\) of the anti-quark. The two diagrams are in figure \ref{fig:LODY}, and we write their amplitudes as  

\begin{align}
    \label{eq:LOme}
    \mathcal{M}(1, 2) = \delta_{ij}\bar{v}(p_2)\gamma^\mu u(p_1)L_\mu(p_\ell, p_{\bar{\ell}})
\end{align}

The function \(L_\mu\) contains the propagator of the photon/Z boson, as well as the lepton vertex. So we just used the electromagnetic vertex of the quarks. 

As is well known, if we have to compute the cross section at NLO, there are two type of corrections we need to compute. Naturally there are the 1 loop diagrams, that are called virtual, because in the sub-process, an off-shell virtual particle is created. Then there are the real diagrams: in those, an extra gluon appears in the final space. This processes are real because in those a real on-shell gluon is created, and in particular the final phase space is different. Real diagrams are tree level: they should be easy, but the integration over the final state gluon gives some trouble. 

Let us discuss in detail. In figure \ref{fig:NLODY} we've drawn the relevant diagrams, both real and virtual. Let's focus on the real. Their amplitude is obtained by eq. \ref{eq:LOme}, by adding the multiplicative factor of the extra color vertex, times the quark propagator, that can take values of \(p_1 - p_4\), or \(-(p_2-p_4)\), if \(p_4\) is the momentum of the emitted gluon. 

\begin{figure}
    \caption{the LO diagrams for DY process}
    \label{fig:LODY}
    \begin{center}
        \begin{tikzpicture}
        \begin{feynman}
            \vertex (a) {\(q\)}; 
            \vertex [below right = of a] (b);
            \vertex [below left= of b] (c) {\(\bar{q}\)}; 
            \vertex [right = of b] (d);  
            \vertex [below right= of d] (e) {\(\ell\)};
            \vertex [above right= of d] (f) {\(\bar{\ell}\)}; 
            \diagram{
                (a) -- [fermion] (b); 
                (b) -- [fermion] (c); 
                (b) -- [boson, edge label = \(\gamma/Z\)] (d); 
                (e) -- [fermion] (d); 
                (d) -- [fermion] (f); 
            }; 
        \end{feynman}
        \end{tikzpicture}
    \end{center}
\end{figure}

\begin{figure}
    \caption{NLO diagrams for DY process}
    \label{fig:NLODY}
    \begin{multicols}{2}
    \begin{center}
        \begin{tikzpicture}
        \begin{feynman}
            \vertex (a) {\(q\)}; 
            \vertex [right = of a] (b);
            \vertex [below = of b] (c) ; 
            \vertex [right = of c] (d);  
            \vertex [right = of b] (g); 
            \vertex [below right= of g] (e) {\(\ell\)};
            \vertex [above right= of g] (f) {\(\bar{\ell}\)}; 
            \vertex [left = of c] (h) {\(\bar{q}\)}; 
            \diagram{
                (a) -- [fermion] (b); 
                (b) -- [fermion] (c); 
                (c) -- [gluon, momentum = \(p_4\)] (d); 
                (e) -- [fermion] (g); 
                (g) -- [fermion] (f); 
                (b) -- [boson, edge label = \(\gamma/Z\)] (g);
                (c) -- [fermion] (h); 
            }; 
        \end{feynman}
        \end{tikzpicture}
    \end{center}
    \columnbreak
    \begin{center}
        \begin{tikzpicture}
        \begin{feynman}
            \vertex (a) {\(q\)}; 
            \vertex [right = of a] (b);
            \vertex [below = of b] (c) ; 
            \vertex [right = of c] (d);  
            \vertex [below right= of d] (e) {\(\ell\)};
            \vertex [above right= of d] (f) {\(\bar{\ell}\)}; 
            \vertex [right = of b] (g); 
            \vertex [left = of c] (h) {\(\bar{q}\)}; 
            \diagram{
                (a) -- [fermion] (b); 
                (b) -- [fermion] (c); 
                (c) -- [boson, edge label = \(\gamma/Z\)] (d); 
                (e) -- [fermion] (d); 
                (d) -- [fermion] (f); 
                (b) -- [gluon, momentum = \(p_4\)] (g);
                (c) -- [fermion] (h); 
            }; 
        \end{feynman}
        \end{tikzpicture}
    \end{center}
    \end{multicols}
\end{figure}


Thus, we write 

\begin{align}
    \label{eq:NLOme}
    \mathcal{M}(1, 2, 4) = \eps^*_\nu(p_4)g_{s, b}(T^a)_{ij}\bar{v}(p_2)\left( \gamma^\mu\frac{\slashed{p_1}-\slashed{p_4}}{(p_1-p_4)^2}\gamma^\nu - \gamma^\nu \frac{ \slashed{p_2}-\slashed{p_4}}{(p_2-p_4)^2}\gamma^\mu\right) u(p_1)L_\mu (p_\ell, p_{\bar{\ell}})
\end{align}

\subsection{The soft limit}

To simplify things, we'd like to use Dirac's equation, that tells us \(\slashed{p_1}u(p_1) = 0 = \bar{v}(p_2)\slashed{p_2}\) (we will always consider massless quarks). In both cases, we anti-commute with the \(\gamma^\nu\) Dirac gamma, 

\begin{align}
    &\slashed{p_1}\gamma^\nu u(p_1) = \left\{ \slashed{p_1}, \gamma^\nu \right\} u(p_1) = 2p_1^\nu u(p_1) \notag \\ 
    &\bar{v}(p_2)\gamma^\nu\slashed{p_2} = 2\bar{v}(p_2)p_2^\nu 
\end{align}

Now consider the "soft" limit, where \(E_4 \to 0\): then we can take \(\slashed{p}_4 = 0\) in the numerator. For the denominator instead we use \((p_1-p_4)^2 = -2p_1 \cdot p_4 = -2E_1 E_4 (1-\cos\theta_{14}) = -2E_1 E_4 \rho_{14}\). In general, we indicate \(\rho_{ij} = (1 - n_i \cdot n_j)\), if \(n_i\) are unit vector that point in the same direction as \(p_i\). Therefore \(\rho_{ij} \in [0, 2], \, \forall i, j\).   

We define \(S_i A := \lim_{E_i \to 0}A\). We find 

\begin{align}
    S_4 \mathcal{M}(1, 2, 4) = g_{s, b}\eps^*_\nu (p_4) (T^a)_{ij} \bar{v}(p_2)\left( -\frac{\gamma^\mu p_1^\nu}{E_1 E_4\rho_{14}}  + \frac{\gamma^\mu p_2^\nu}{E_2 E_4\rho_{24}}\right)u(p_1)L_\mu (p_\ell, p_{\bar{\ell}})
\end{align}

squaring this, we can write 

\begin{align}
    \label{eq:softDYfactorization}
    S_4 |\mathcal{M}(1, 2,4)|^2 &= -\frac{g_{s, b}^2 (N_c^2-1)}{2N_c E_4^2} \left(-\frac{p_1^\nu}{E_1 \rho_{14}}  +\frac{p_2^\nu}{E_2\rho_{24}}\right)^2|\mathcal{M}(1, 2)|^2 \notag \\ 
    &= 2\frac{g_{s, b}^2 C_F}{E_4^2}\left( \frac{p_1\cdot p_2}{E_1E_2 \rho_{14}\rho_{24}} \right)|\mathcal{M}(1, 2)|^2 = 2\frac{g_{s, b}^2 C_F}{E_4^2}\left( \frac{\rho_{12}}{\rho_{14}\rho_{24}} \right)|\mathcal{M}(1, 2)|^2
\end{align}

In the first line, the factor \(N_c^2-1\) comes from \(\text{Tr}(T^a T^a) = \frac{1}{2}\delta^{aa} = \frac{N_c^2 -1}{2}\), and the factor \(N_c^-1\) is the color factor of \(|\mathcal{M}(1, 2)|^2\). Thus we see \(C_F = \frac{N_c^2 -1}{2N_c}\) appearing naturally. This is the known factorization of the matrix element in the soft limit, and the factor before \(|\mathcal{M}(1, 2)|^2\) is called \emph{eikonal} factor. In can be shown that this is happens universally, not only for DY process.

Now that we have properly obtained the eikonal factor, we see that its integral over the gluon finale phase space is not finite. Of course this is a problem: we cannot, for example, put the integral in a MC program and get a sensibile result. 

This is not the only singularity of the amplitude squared. In the soft limit, the denominator cointains the factors \(\rho_{14}\rho_{24}\propto (1-\cos\theta_{14})(1-\cos\theta_{24})\) that go to zero as \(\theta^{2}\) as \(\theta \to 0\). Thus also the angular integral is divergent. 

In this section we focused on a particular process (DY), but we used a rather general approach. Indeed, one can show that the factorization of the amplitude in the soft limit is universal, and given by the following formula (Catani, 2000): 

\begin{align}
    \label{eq:Catanisoftfact}
    \lim_{E_q \to 0 }|\mathcal{M}(p_1, \ldots, p_n; q)|^2 = -4\pi\alpha_s \mu^{2\eps}\sum_{ij}S_{ij}\vec{T}_i \cdot \vec{T}_j |\mathcal{M}(p_1, \ldots, p_n)|^2, \quad S_{ij} = \frac{p_i\cdot p_j}{(p_i\cdot q)(p_j \cdot q)}
\end{align}

Let's show that this reproduces eq \ref{eq:softDYfactorization}. Remember that \(g_{s, b}^2 = 4\pi\alpha_s\). For now, let's not worry about UV divergences and the running of \(\alpha_s\). The new part is computing the color products. In general, one can compute them using \((\vec{T}_q)^2 = (\vec{T}_{\bar{q}})^2 = C_F\), as well as color conservation. The sum runs over the distinct indexes \(i, j\) and is symmetric in the exchange. So we find 

\begin{align}
    -2g_{s, b}^2 \frac{p_i\cdot p_j}{(p_i\cdot q)(p_j\cdot q)}(\vec{T}_q \cdot \vec{T}_{\bar{q}})|\mathcal{M}(1, 2)|^2 
\end{align}

Using color conservation, 

\begin{align}
    \vec{T}_q + \vec{T}_{\bar{q}} = 0 \implies (\vec{T}_q + \vec{T}_{\bar{q}})^2 = 2C_F +2\vec{T}_q \cdot \vec{T}_{\bar{q}} = 0 \implies \vec{T}_q \cdot \vec{T}_{\bar{q}} = -C_F 
\end{align}

So we recover our previous expression. 

\subsection{The collinear limit}

That the angular integral is divergent in the collinear limit (i.e. when one of the \(\theta\)s goes to zero) is true also when the gluon emmitted is \emph{hard}. To see this we go back to the amplitude, eq. \ref{eq:NLOme}. 

\section{Dealing with IR divergences: pdfs renormalization}

If the result is infinite, at some points in our computation we must have done a mistake. Actually, the mistakes are two: 

\begin{enumerate}
    \item We did not consider the virtual processes that contribute to the same order in \(\alpha_s\). 
    \item We computed a partonic cross section: alone, it is not a physical observable as we can't have free partons in the initial state. 
\end{enumerate}

To obtain a sensible result, first we need to add up the virtual corrections, and then we need to integrate the cross section using the partonic distribution functions (pdfs). Thus, we can write the pdfs as a sum of a bare pdf and a counterterm, that eliminates the remaining poles. 

The \emph{hadronic} (thus observable) cross section is given from \(F_{LM}(1, 2)\) as 

\begin{align}
    \label{eq:hadcrosssec}
    \int_0^1 dx_1 dx_2 f_q^{bare}(x_1)f_q^{bare}(x_2)\frac{1}{x_1P_1 x_2P_2}F_{LM}(x_1P_1, x_2P_2)
\end{align}

where \(P_1, \, P_2\) are the momenta of the incoming hadrons, and we recognize the flux factor \((2s)^{-1} = (x_1P_1 x_2P_2)^{-1}\). 
 
As done for the couplings or masses, the counterterms for the pdfs can be computed as a perturbative series in \(\alpha_s\). In the \(\bar{\text{MS}}\) scheme, one writes the bare pdfs in terms of the remormalized using 

\begin{align}
    f_i^{bare}(x) = [\delta_{ij}\delta(1-y_2)+\frac{\alpha_s(\mu^2)}{2\pi}\frac{1}{\eps}\hat{P}^{(0)}_{ij}(y_1) + o(\alpha_s^2)] \otimes f_j^{ren}(y_2)
\end{align}

where \(\otimes\) is the convolution. \(\hat{P}^{(0)}_{ij}\) is instead the Altarelli-Parisi kernel function. Inserting this in equation \ref{eq:hadcrosssec}, considering for now only the terms in \(x_1\), one gets 

\begin{align}
    &\int_0^1 dx_1 f_q^{bare}(x_1)\frac{F_{LM}(x_1P_1, \ldots)}{x_1} = \notag \\ 
    &\int_0^1 dx_1 f_q^{ren}(x_1)\frac{F_{LM}(x_1P_1, \ldots)}{x_1} + \int_0^1 dx_1 \frac{\alpha_s(\mu^2)}{2\pi\eps}[\hat{P}^{(0)}_{qq}\otimes f_q^{ren}](x_1) \frac{F_{LM}(x_1P_1, \ldots)}{x_1} 
\end{align}

We recognize in the second term a counterterm of order \(\alpha_s\). Writing explicitly the convolution, and integrating over \(x_1\) using the delta function, we obtain: 

\begin{align}
    &\int_0^1 dx_1 dy_1 dy_2 \frac{\alpha_s(\mu^2)}{2\pi x_1\eps}\hat{P}_{qq}(y_1)f_q^{ren}(y_2)F_{LM}(x_1P_1, \ldots)\delta(x_1- y_1y_2) =\notag \\ 
    &\int_0^1 dy_1 dy_2 \frac{\alpha_s(\mu^2)}{2\pi y_1 y_2\eps}\hat{P}^{(0)}_{qq}(y_1)f^{ren}_{q}(y_2)F_{LM}(y_1y_2 P_1, \ldots) \notag  = \\ 
    &\int_0^1 dx_1 \frac{f^{ren}_q (x_1)}{x_1}\int_0^1 dz \frac{\alpha_s(\mu^2)}{2\pi z} \frac{1}{\eps}\hat{P}^{(0)}_{qq}(z)F_{LM}(zx_1 P_1, \ldots)
\end{align}

We recognize therefore a \emph{partonic} counterterm, given by 

\begin{align}
    d\sigma^{CV, 1} = \frac{\alpha_s(\mu^2)}{2\pi\eps}\int_0^1 dz \hat{P}_{qq}(z)\frac{F_{LM}(z\cdot1, 2)}{z}
\end{align}

Adding the contribution to the other renormalized pdf (\(f^{bare}_{q}(x_2)\)), we get the form of the explicit counterterm: 

\begin{align}
    \label{eq:pdfpartoniccounterterm}
    d\sigma^{CV} = \frac{\alpha_s(\mu^2)}{2\pi\eps}\int_0^1 dz \hat{P}_{qq}^{(0)}(z)\left[ \frac{F_{LM}(z\cdot1, 2)}{z} + \frac{F_{LM}(1, z\cdot 2)}{z} \right]
\end{align}

\section{The subtraction scheme for \(q\bar{q} \to V + g\)}

As mentioned before, if an integral is divergent, we cannot integrate it numerically. A possibile strategy for dealing with this problem is defining some operators that extract the singularities. Then one can try to compute analitically the integral in the particular limit, and use a numerical algorithm (e.g. Monte Carlo integrator) to compute the regular integral.   

To see how we can do this in practice, first we introduce some useful notation. We also switch to the more general case where the process is \(q\bar{q} \to V + g\), with \(V\) colorless particles. We denote the contribution of the real diagrams of \(d\sigma^{NLO}\) as 

\begin{align}
    d\sigma^R = \frac{1}{2s}\int [dg_4] F_{LM}(1, 2, 4)
\end{align}

where \((2s)^{-1}\) is the flux factor, \([dg_4] = \frac{d^{d-1}p_4\theta(E_{max}-E_4)}{(2\pi)^d 2E_4}\) and 

\begin{align}
    F_{LM}(1, 2, 4) = d\text{Lips}_V |\mathcal{M}(1, 2,4)|^2 \mathcal{F}_{kin}(1, 2, 4)
\end{align}

where \(d\text{Lips}_V\) is the lorentz invariant phase space of the particles \(V\), and \(\mathcal{F}_{kin}(1, 2, 4)\) is a IR safe observable. 

Now notice that the soft limit is not the only point in the phase space that is not integrable. There is also the so called \emph{collinear} singularity, that arises if \(\theta_{14}\) or \(\theta_{24}\) vanishes (or, equivalently, \(\rho_{12}\) or \(\rho_{14}\) go to 0). To be able to integrate the matrix element numerically we must therefor subtract both the soft and the collinear divergences. To this aim we also define the collinear operator, 

\begin{align}
    C_{ij}A = \lim_{\rho_{ij} \to 0}A 
\end{align}

The isoleted divergences of \(F_{LM}(1, 2, 4)\) are given by \(C_{41}F_{LM}(1, 2, 4), \, C_{42}F_{LM}(1, 2, 4), \, S_4F_{LM}(1, 2, 4)\). We thus re-write \(d\sigma^R\) as (using the convention where \(\int [dg_4] \ldots = <\ldots> \))

\begin{align}
    d\sigma^R &= \frac{1}{2s}<(C_{41}+C_{42})F_{LM}(1, 2, 4)> + \frac{1}{2s}<(1 - (C_{41}+C_{42}))F_{LM}(1, 2, 4)> \notag \\ 
    &= \frac{1}{2s}(<(C_{41} + C_{42})F_{LM}(1, 2, 4)> + <(I - C_{41} - C_{42})S_4F_{LM}(1, 2, 4)> + <\hat{O}_{NLO}F_{LM}(1, 2, 4)>) 
\end{align}

where we defined \(\hat{O}_{NLO} = (I-C{41}- C_{42})(I-S_4)\). The term \(\propto \hat{O}_{NLO}\) is now free from any singularity by costruction, and can be numerically computed safely. Hence we must compute only the two extra terms. 

\clm{}{}{

    The term \( <(I - C_{41} - C_{42})S_4F_{LM}(1, 2, 4)>\) is zero 
}

To see this, we use the expresion for \(S_4 F_{LM}(1, 2, 4)\) that we derived earlier (eq. \ref{eq:softfactorization}). In the general case, the formula can be written as 

\begin{align}
    S_4 F_{LM}(1, 2, 4) = \frac{2g_{s,b}^2 C_F}{E_4^2}\frac{\rho_{12}}{\rho_{14}\rho_{24}}F_{LM}(1, 2)
\end{align}

We have to compute the action of \(C_{14}\) as well as \(C_{24}\) on this. Of course the limit affects only the fraction \(\frac{\rho_{12}}{\rho_{14}\rho_{24}}\). In the CM frame, \(\rho_{12} = 2\). Also, if \(p_4\) is allined with \(p_1\), it must be anti-parallel to \(p_2\). We conclude that 

\begin{align}
    C_{14}\frac{\rho_{12}}{\rho_{14}\rho_{24}} = \frac{1}{\rho_{14}}, \quad C_{24}\frac{\rho_{12}}{\rho_{14}\rho_{24}} = \frac{1}{\rho_{24}} \implies (C_{14} + C_{24}) \frac{\rho_{12}}{\rho_{14}\rho_{24}} = \frac{1}{\rho_{14}} + \frac{1}{\rho_{24}} = \frac{\rho_{14} + \rho_{24}}{\rho_{14}\rho_{24}}
\end{align}

Again since we're in the CM frame, \(\rho_{14} + \rho_{24} = 2 = \rho_{12}\). We conclude that \(C_{14} + C_{42}\) leaves \(SF_{LM}(1, 2, 4)\) unchanged, and so \((I - C_{41} - C_{42})S_4F_{LM}(1, 2, 4) = 0\). 

The only term that we have to deal with is therefore 

\begin{align}
    <(C_{41} + C_{42})F_{LM}(1,2,4)>
\end{align}

When we derived the collinear limit factorization for the DY process (eq. ), we mentioned that it was completely universal. We can safely write than 

\begin{align}
    \label{eq:collimNLO}
    C_{41}F_{LM}(1, 2, 4) = \frac{g^2_{s,b}}{E_4^2\rho_{14}}\frac{1-z}{z}P_{qq}(z)F_{LM}(z\cdot1, 2)
\end{align}

The splitting function is defined as before, and \(z = 1- E_4/E_1\). We also simplified the denominator using \((p_1-p_4)^2 = -2p_1\cdot p_4 = -2E_1E_4\rho_{14} = -2E_4^2\rho_{14}\frac{E_1}{E_4} = -2\frac{E_4^2\rho_{14}}{1-z}\). 

Problems start to arise when we consider this term integrated over the gluon final phase space. We need to evaluate the integral in dimensional regularization, so we can cancel it with the UV divergences of the virutual NLO corrections. We write then 

\begin{align}
    [dg_4] = \frac{E_4^{d-3}dE_4 d\cos\theta (1-\cos^2\theta)^{d/2-2} d\Omega_{d-2}}{2(2\pi)^{d-1}}
\end{align}

since \(E_4 = |p_4|\). Our integrand (eq. \ref{eq:collimNLO}) depends on \(E_4\) through \(z\), and the angular dependence is only in \(\rho_{14}^{-1}\). Our next step is therefore to compute the angular integral, and to write the energy integral in terms of \(z\).  

Let's start from the angular integral, that simply reads 

\begin{align}
    \int d\Omega_{d-2} d\cos\theta \frac{(1-\cos^2\theta)^2}{1-\cos\theta} = \int d\Omega_{d-2}\int d\cos\theta (1-\cos\theta)^{d/2-3}(1+\cos\theta)^{d/2-2}
\end{align}

The last integral strongly resembles the euler \(\beta\) integral definition, 

\begin{align}
    B(\alpha, \gamma) = \int_0^1 dx x^{\alpha-1}(1-x)^{\gamma-1} = \beta(\alpha, \gamma) = \frac{\Gamma(\alpha)\Gamma(\gamma)}{\Gamma(\alpha+\gamma)}
\end{align}

So we use the substitution \(x = \frac{1}{2}(1-\cos\theta)\): 

\begin{align}
    \int d\Omega_{d-2} 2^{d-4}\int_0^1 dx (x)^{d/2-3}(1-x)^{d/2-2} &= \int d\Omega_{d-2}2^{d-4}\frac{\Gamma(d/2-2)\Gamma(d/2-1)}{\Gamma(d-3)} \notag \\ 
    &= -\frac{1}{\eps 4^{\eps}}\frac{\Gamma^2(1-\eps)}{\Gamma(1-2\eps)}\int d\Omega_{d-2}
\end{align}

We now use \([\alpha_s] = \int d\Omega_{d-2}\frac{g_{s, b}^2\mu_0^{2\eps}}{2(2\pi)^{d-1}}\) (the factor \(\mu_0^{2\eps}\) is needed for dimensional reasons, since from now on we will take \(d = 4-2\eps\)), and the full integral reads

\begin{align}
    \braket{C_{41}F_{LM}(1,2, 4)} = -\frac{[\alpha_s]}{\eps 4^{1+\eps}}\frac{\Gamma^2(1-\eps)}{\Gamma(1-2\eps)}\int dE_4  E_4^{d-5}\theta(E_{max} - E)\frac{1-z}{z}P_{qq}(z)F_{LM}(z\cdot 1, 2)
\end{align}

Then, we use \(z = 1-E_4/E_1 \implies E_4 = (1-z)E_1, \, \,  dE_4 = -E_1 dz \), 

\begin{align}
    \label{eq:collzintegral}
    \braket{C_{41}F_{LM}(1,2, 4)} &= -\frac{[\alpha_s]}{\eps 4^{1+\eps}}\frac{\Gamma^2(1-\eps)}{\Gamma(1-2\eps)}\int_{z_{\min}}^1 dz E_1  [(1-z)E_1]^{-(1+2\eps)}\frac{1-z}{z}P_{qq}(z)F_{LM}(z\cdot 1, 2) \notag \\ 
    &= -\frac{[\alpha_s]}{\eps}\frac{\Gamma^2(1-\eps)}{\Gamma(1-2\eps)}(2E_1)^{-2\eps}\int_{z_{\min}}^1 \frac{dz}{(1-z)^{2\eps}}P_{qq}(z) \frac{F_{LM}(z\cdot 1, 2)}{z}
\end{align}

Now we need to write to write the \(z\) integral in such a way that the combination with the virtual corrections is straightforward, i.e. the pole structure is completely explicit. This is still clearly not the case, as in the limit \(z\to 1\) the integral is divergent (or at least would be divergent in the limit \(\eps \to 0\)): this is of course the soft limit. We write \(F_{LM}(1,2,4)/z = G(z)\), and we rewrite \(P_{qq}(z)\) in the following way:  

\begin{align}
    P_{qq}(z) &= \frac{C_F}{1-z} + \frac{C_F z^2}{1-z} - C_F\eps(1-z) = \frac{2C_F}{1-z} + \frac{C_F(z^2 -1)}{1-z} -C_F\eps(1-z) \notag \\ 
    &= \frac{2C_F}{1-z} + P_{qq}^{reg}(z)
\end{align}

In doing so we explicitly separated \(P_{qq}\) in a a piece that is singular in the limit \(z\to1\), \(\frac{2C_F}{1-z}\), and one piece that is regular \(P^{rel}_{qq} = -C_F((z+1) + \eps(1-z))\). 

Consider then the integral in \(z\): (we can safely extend the integral to the range \([0, 1]\) as \(F_{LM}\) contains the final state phase space that is zero if there is not enough energy to produce the final state particles, i.e. if \(z<z_{\min}\))

\begin{align}
    I_z &= \int_0^1 dz \left[ \frac{2C_F}{(1-z)^{1+2\eps}}+(1-z)^{-2\eps}P_{qq}^{reg}(z) \right]G(z) = \notag \\ 
    &= \int_0^1 dz \left[ \frac{2C_F(G(z)-G(1))}{(1-z)^{1+2\eps}}+(1-z)^{-2\eps}P_{qq}^{reg}(z)G(z) \right] + X
\end{align}

The last expression holds only if 

\begin{align}
    X = \int_0^1 dz \frac{2C_FG(1)}{(1-z)^{1+2\eps}} = \frac{2C_FG(1)}{-2\eps} = -\frac{C_F G(1)}{\eps}
\end{align}

We notice than that we've taken the divergent part out of the integral. The first fraction is not divergent anymore as also the numerator now goes to zero as \(z\to1\). More explicitly, we write \(\frac{G(z)-G(1)}{(1-z)^{1+2\eps}}\)as a power series: 

\begin{align}
    \frac{G(z)-G(1)}{(1-z)^{1+2\eps}} &= (1-z)^{-2\eps}\left( \frac{G(z)-G(1)}{1-z} \right) = e^{-2\eps\log(1-z)}\left( \frac{G(z)-G(1)}{1-z} \right) = \notag \\ 
    &= \sum_{n=0}^{\infty} \frac{(-1)^n (2\eps)^n \log^n(1-z)}{n!}\left( \frac{G(z)-G(1)}{1-z} \right) 
\end{align}

The first part of the integral then becomes  

\begin{align}
    2C_F \sum_{n=0}^{\infty}\frac{(-2\eps)^n}{n!}\int_0^1 \frac{\log^n(1-z)}{(1-z)}(G(z)-G(1)) = 2C_F \sum_{n=0}^{\infty}\frac{(-2\eps)^n}{n!}\int_0^1 \mathcal{D}_n(z) G(z)  
\end{align}

The factors \(\mathcal{D}_n\) are defined by the last equality and are called \emph{plus-distributions}.

In equation \ref{eq:collzintegral} the \(z\) integral is moltiplied by a pole \(1/\eps\). Therefore we must expand the integral up to \(o(\eps)\). The whole integral now reads

\begin{align}
    I_z = -\frac{C_F F_{LM}(1, 2)}{\eps} + \underbrace{2C_F \sum_{n=0}^{\infty}\frac{(-2\eps)^n}{n!}\int_0^1 \mathcal{D}_n(z) G(z)}_{=I_z^1} + \underbrace{\int_0^1 dz (1-z)^{-2\eps}P_{qq}^{reg}(z)}_{=I_z^2}
\end{align}

The first term is already explicit, so we just need to expand \(I_z^1, I_z^2\): 

\begin{align}
    &I_z^1 = 2C_F \left( \int_0^1 dz D_0(z) G(z)  -2\eps \int_0^1 dz D_1(z)G(z)\right) \notag \\ 
    &I_z^2 = \int_0^1 (1-z)^{-2\eps}P_{qq}^{reg}(z) = -C_F\int_0^1 dz e^{-2\eps\log(1-z)}(z+1 - \eps(1-z)) = \notag \\ 
    &\quad = -C_F\int_0^1 [(z+1) -2\eps\log(1-z)(z+1) + \eps(1-z)G(z) 
\end{align}

Combining everything, we find: 

\begin{align}
    \label{eq:Izres}
    I_z = -\frac{C_F F_{LM}(1, 2)}{\eps} + C_F \int_0^1 dz (2\mathcal{D}_0 - (1+z)) + \eps[2\log(1-z)(z+1) - (1-z) -4\mathcal{D}_1(z)]G(z) + o(\eps^2)
\end{align}

Our result is then 

\begin{align}
    \label{eq:colintegralres}
    \braket{C_{41}F_{LM}(1, 2, 4)} = -\frac{[\alpha_s]}{\eps}\frac{\Gamma^2(1-\eps)}{\Gamma(1-2\eps)}s^{-\eps}\left[-C_F \left(\frac{1}{\eps} +\frac{3}{2} \right)\braket{F_{LM}(1, 2)} + \int_0^1 \hat{P}_{qq, R}(z)\braket{\frac{F(z\cdot 1, 2)}{z}}\right]    
\end{align}

where \(\hat{P}_{qq, R}\) is a modified splitting function, explicitly given by 

\begin{align}
    &\hat{P}_{qq, R} = \hat{P}_{qq}^{(0)} + \eps \hat{P}_{qq}^{(\eps)} + o(\eps^2),  \notag \\ 
    &\hat{P}_{qq}^{(0)} = C_F\left(2\mathcal{D}_0 -(1+z) + \frac{3}{2}\delta(1-z)\right) \notag \\ 
    &\hat{P}_{qq}^{(\eps)} = C_F(2\log(1-z)(z+1)- (1-z) - 4\mathcal{D}_1(z))
\end{align}

the expression for \(\hat{P}^{(0)}_{qq}\) matches the LO Altarelli-Parisi splitting kernel. Notice that to obtain eq. \ref{eq:colintegralres} we just used our result for the z integral (eq. \ref{eq:Izres}) into \ref{eq:collzintegral}, adding and subtracting the factor \(-3/2\). 

Having obtained this expression for the first collinear limit, the second one will be analogous. We can then write the real cross section, in the explicit form that we wanted: 

\begin{align}
    2sd\sigma^{R} = &\braket{\hat{O}_{NLO}F_{LM}(1, 2, 4)} + 2C_F[\alpha_s]\frac{\Gamma^2(1-\eps)}{\Gamma(1-2\eps)}s^{-\eps}\left[ \left(\frac{1}{\eps^2} +\frac{3}{2\eps} \right)\braket{F_{LM}(1, 2)}\right] \notag \\ 
    &-\frac{[\alpha_s]}{\eps}\frac{\Gamma^2(1-\eps)}{\Gamma(1-2\eps)}s^{-\eps}\int_0^1 \hat{P}_{qq, R}(z)\left(\braket{\frac{F(z\cdot 1, 2)}{z}} + \braket{\frac{F(1, z\cdot 2)}{z}}  \right)
\end{align}

We make some comments on this result. The first part is integrable, and can be integrated numerically. The second part is the most divergent (\(\propto \eps^{-2}\)), and multiplies directly the LO factor, \(\braket{F_{LM}(1, 2)}\). Those poles are what we expect to be canceled together with the poles form the virtual diagrams. Lastly, the last term, proportional to \(\eps^{-1}\), is proportional to the LO amplitude with rescaled energy. Such factors \emph{do not} arise when computing UV divergences, and therefore must be canceled with the collinear counterterms. 

Let's now add the UV poles that come from the virtual diagrams. We use Catani's representation of renormalized one loop amplitudes: 

\begin{align}
    2sd\sigma^{V} = -2[\alpha_s]\cos(\eps \pi)\left( \frac{C_F}{\eps^2}+\frac{3C_F}{2\eps} \right)s^{-\eps}\braket{F_{LM}(1, 2)} + \braket{F_{LV}^{fin}(1, 2)}
\end{align}

The sum of the two contributes is thus 

\begin{align}
    2sd\sigma^{V+R} = &\braket{\hat{O}_{NLO}F_{LM}(1, 2, 4)} + \braket{F_{LV}^{fin}(1, 2)} + \notag \\
    &2[\alpha_s]s^{-\eps}\left( \frac{\Gamma^2(1-\eps)}{\Gamma(1-2\eps)} -\cos(\pi\eps)\right)\left( \frac{C_F}{\eps^2} +\frac{3C_F}{2\eps}\right)\braket{F_{LM}(1, 2)} + \notag \\ 
    &- \frac{[\alpha_s]}{\eps}\frac{\Gamma^2(1-\eps)}{\Gamma(1-2\eps)}s^{-\eps}\int_0^1 dz \hat{P}_{qq, R}\left[\frac{F_{LM}(z\cdot1, 2)}{z} + \frac{F_{LM}(1, z\cdot 2)}{z}  \right]
\end{align}

Expanding the function \(\mathcal{S}(1, 2) = s^{-\eps}\left( \frac{\Gamma^2(1-\eps)}{\Gamma(1-2\eps)} -\cos(\pi\eps)\right)\left( \frac{C_F}{\eps^2} +\frac{3C_F}{2\eps}\right)\) up to \(o(1)\), one finds that all the poles are calceled, leaving 

\begin{align}
    2sd\sigma^{V+R} = &\braket{\hat{O}_{NLO}F_{LM}(1, 2, 4)} + \braket{F_{LV}^{fin}(1, 2)} + \frac{2\pi^2C_F}{3}[\alpha_s]s^{-\eps}\braket{F_{LM}(1, 2)} \notag \\
    &- \frac{[\alpha_s]}{\eps}\frac{\Gamma^2(1-\eps)}{\Gamma(1-2\eps)}s^{-\eps}\int_0^1 dz \hat{P}_{qq, R}\left[\frac{F_{LM}(z\cdot1, 2)}{z} + \frac{F_{LM}(1, z\cdot 2)}{z}  \right]
\end{align}

All is left is adding the pdfs counterterm. In the \(\bar{\text{MS}}\) scheme, this means adding the factor in eq. \ref{eq:pdfpartoniccounterterm} to the partonic cross section. We therefore find 

\begin{align}
    2sd\sigma^{NLO} = &\braket{\hat{O}_{NLO}F_{LM}(1, 2, 4)} + \braket{F_{LV}^{fin}(1, 2)} + \frac{2\pi^2C_F}{3}[\alpha_s]s^{-\eps}\braket{F_{LM}(1, 2)} + \notag \\ 
    & \frac{\alpha_s(\mu^2)}{2\pi\eps}\left(1-\frac{\mu^{2\eps}s^{-\eps} e^{\eps\gamma_E}\Gamma(1-\eps)}{\Gamma(1-2\eps)}\right)\int_0^1 dz \hat{P}^{0}_{qq}(z) \left[\frac{F_{LM}(z\cdot1, 2)}{z} + \frac{F_{LM}(1, z\cdot 2)}{z}  \right] \notag \\ 
    &- [\alpha_s]\frac{\Gamma^2(1-\eps)}{\Gamma(1-2\eps)}s^{-\eps}\int_0^1 dz \hat{P}^{(\eps)}_{qq, R}\left[\frac{F_{LM}(z\cdot1, 2)}{z} + \frac{F_{LM}(1, z\cdot 2)}{z}  \right]
\end{align}

Notice that the third line is free from any pole. Expanding the \(\eps\) dependent terms in the second line up to \(o(\eps)\) we get:  

\begin{align}
    \left( 1 - \frac{\mu^{2\eps}s^{-\eps}\Gamma(1-\eps)}{\Gamma(1-2\eps)} \right) &= 1 - e^{\eps(\log(\mu^2/s) +\gamma_E)}\frac{1+\eps\gamma_E}{1+2\eps\gamma_E} \notag \\ 
    &= 1 - (1-\eps(\log(\mu^2/s)) + \eps\gamma_E)(1-\eps\gamma_E) + o(\eps^2) = \eps\log\left( \frac{s}{\mu^2} \right) + o(\eps^2)
\end{align}

The poles cancel completely! The final result is then


\begin{align}
    2sd\sigma^{NLO} = &\braket{\hat{O}_{NLO}F_{LM}(1, 2, 4)} + \braket{F_{LV}^{fin}(1, 2)} + \frac{2\pi^2C_F}{3}\frac{\alpha_s(\mu^2)}{2\pi}\braket{F_{LM}(1, 2)} + \notag \\ 
    & \frac{\alpha_s(\mu^2)}{2\pi}\int_0^1 dz \left[\hat{P}^{(0)}_{qq}(z)\log\left( \frac{s}{\mu^2} \right) - \hat{P}^{(\eps)}_{qq, R}\right] \left[\frac{F_{LM}(z\cdot1, 2)}{z} + \frac{F_{LM}(1, z\cdot 2)}{z}  \right]
\end{align}

where we also performed the limit \(\eps \to 0\). 

\section{A first generalization: the case \(q\bar{q} \to X + Ng\)}

For the hadronic process \(pp \to X + N\) jets , the full cross section is given in terms of the partonic cross sections as 

\begin{align}
    \label{eq:qqXNgcrosssec}
    d\sigma = \sum_{a, b} \int dx_1 dx_2 f_a(x_1, \mu_F)f_b(x_2, \mu_F)d\hat{\sigma}_{ab}(x_1, x_2, \mu_R, \mu_F; O)
\end{align}

Where \(\mu_{R}\) is the renormalization scale, \(\mu_F\) the factorization scale. \(O\), instead, is an appropriate observable, that also provides an IR safe definitions of the \(N\) jets. This means that the tree level amplitude is safe from IR divergences (this is good news of course, because at tree level we wouldn't be able to cancel those divergences adding up the IR divergences of virtual processes, as well as the pdfs counterterms). For the moment, we're interested in \(d\hat{\sigma}_{ab}^{NLO}\), that is given (exactly as the last section) by the sum of 3 distinct terms: 

\begin{align}
    d\hat{\sigma}^{NLO}_{ab} = d\hat{\sigma}^{R} + d\hat{\sigma}^{V} + d\hat{\sigma}^{pdf}
\end{align}

the three terms are, in order, the real, virtual and pdf counterterm cross section. 

Let's start with the real emission cross section. This process is \(1_a + 1_b \to X + (N+1)g\). It is again a tree level amplitude, with an extra gluon in the final phase space. 

An important point is that since the observable \(O\) makes sure that at least \(N\) partons are resolved, only one gluon can become soft or collinear. We must therefore be able to identify this unresolved parton. We introduce than the damping factors \(\Delta^{(i)}\), \(i \in 3, \ldots, N_p+1\), (\(N_p\) is the number of partons involved in the process, so for \(pp \to X + N g\), \(N_p = N+2\) and we always leave the indexes \(1, 2\) for the initial state partons) that are defined in such way that they go to zero for every \(i\), except if the parton \(i\) is unresolved. They also provide a partition of unity, 

\begin{align}
    \sum_{i=3}^{N_p+1}\Delta^{(i)} = 1
\end{align}

Following the same scheme of the previous section (and using the same notation), we want to compute 

\begin{align}
    \braket{F_{LM}(1, 2; 3, \ldots, N_p +1, X)} = \sum_{i=3}^{N_p+1}\braket{\Delta^{(i)}F_{LM}(1, 2; 3, \ldots, N_p +1, X)}
\end{align}

In this last equation, \(F_{LM}(\ldots)\) is completely symmetric in any exchange of the \(N +1\) external gluons. Then in each term of the sum, we can re-label the index, in such a way that all the damping functions are identical, and we obtain 

\begin{align}
    \sum_{i=3}^{N_p+1}\braket{\Delta^{(i)}F_{LM}(1, 2; 3, \ldots, N_p +1, X)} = \braket{(N_p-1)\Delta^{\mathfrak{m}}F_{LM}(\mathfrak{m})}
\end{align}

The factor \(N_p-1\) multiplies the \(1/(N_{p}-1)!\), turning it into \(1/(N_p-2)! = 1/N!\), where \(N\) are the number of resolved jets. This is therefore the same symmetry factor that sits in the \(F_{LM}\) of the virtual processes. To the sake of IR poles combinations we can omit in both functions the explicit symmetry factor.   

As we did in the previous section, we now want to extract the divergences, using the operators \(S_i, \, C_{ij}\). First, we list the action of those operators on the damping functions, 

\begin{align}
    S_{\mathfrak{m}} \Delta^{(\mathfrak{m})} = 1, \quad C_{a\mathfrak{m}}\Delta^{(\mathfrak{m})} = 1, \quad C_{i\mathfrak{m}}\delta^{(\mathfrak{m})} = \frac{E_{\mathfrak{m}}}{E_{\mathfrak{m}} + E_i} =: z_{\mathfrak{m}, i}
\end{align}

where \(a = 1, 2\) and \(i = 3, \ldots, N_p +1\), as before. We start by subtracting the soft, 

\begin{align}
    \braket{F_{LM}(\mathfrak{m})} = \braket{S_{\mathfrak{m}}F_{LM}(\mathfrak{m})} + \braket{\bar{S}_{\mathfrak{m}}F_{LM}(\mathfrak{m})}
\end{align}

For any operator \(A\), we denote \(\bar{A} = I -A\). In the previous section, we only had to consider the collinear limit, because in that case we found that acting with the collinear operators on the soft limit leaves the expression invariant. This is not the case anymore, as the gluon can go collinear to both initial state and final state resolved partons.  

Therefore we must already write \(\braket{S_{\mathfrak{m}}F_{LM}(\mathfrak{m})}\), expliciting the \(\eps\) pole structure. We use the universal form of the soft gluon limit (eq. \ref{eq:Catanisoftfact}), and we integrate it over the gluon final phase space: 

\begin{align}
    \braket{S_{\mathfrak{m}}F_{LM}(\mathfrak{m})} = -\frac{g_{s, b}^2\mu^{2\eps}}{2(2\pi)^{d-1}}\int dE_{\mathfrak{m}} E_{\mathfrak{m}}^{d-5}\theta(E_{max} - E)d^{d-2}\Omega \sum_{(ij)}\frac{\rho_{ij}}{\rho_{i\mathfrak{m}}\rho_{j\mathfrak{m}}}(\vec{T}_i \cdot \vec{T}_j)\braket{F_{LM}}
\end{align}

From now one we use \(d = 4-2\eps\), and we inserted the scale \(\mu^{2\eps}\). The energy integral is straightforward,

\begin{align}
    \braket{S_{\mathfrak{m}}F_{LM}(\mathfrak{m})} = -\frac{g_{s, b}^2}{2(2\pi)^{d-1}}\left( \frac{E_{\max}}{\mu}\right)^{-2\eps}\int d^{d-2}\Omega \sum_{(ij)}\frac{\rho_{ij}}{\rho_{i\mathfrak{m}}\rho_{j\mathfrak{m}}}(\vec{T}_i \cdot \vec{T}_j)\braket{F_{LM}}
\end{align}

The angular integral is slightly more complicated, and can be expressed in terms of Gauss's hypergeometric function, \(\, _2F_1(a, b, c,d)\). The whole integral reads 

\begin{align}
    \label{eq:intsoftlimNLO}
    \braket{S_{\mathfrak{m}}F_{LM}(\mathfrak{m})} = -[\alpha_s]\frac{(2E_{\max}/\mu)^{-2\eps}}{\eps^2}\sum_{(ij)}^{N_p}\braket{\eta_{ij}^{-\eps}K_{ij}(\vec{T}_i \cdot \vec{T}_j)F_{LM}}
\end{align}

We now proceed in substituting the singular limits, and define a partition of the phase space, that has the following properties: 

\begin{align}
    \sum_{i=1}^{N_p}\omega^{\mathfrak{m}i} = 1, \quad C_{j\mathfrak{m}}\omega^{i\mathfrak{m}} = \delta_j^i
\end{align}

We find 

\begin{align}
    \label{eq:intFLMm}
    \braket{F_{LM}(\mathfrak{m})} = &\braket{S_{\mathfrak{m}}F_{LM}(\mathfrak{m})} + \sum_i \braket{\bar{S}_{\mathfrak{m}}C_{\mathfrak{m} i}F_{LM}(\mathfrak{m})} + \notag \\ 
    &+\sum_i \braket{\bar{S}_{\mathfrak{m}}\bar{C}_{\mathfrak{m} i} \omega^{i\mathfrak{m}} F_{LM}(\mathfrak{m})}
\end{align}

Now the last line has all the singularities removed. The first one we have computed, with all the \(\eps\) poles fully explicit. Thus we need to compute the second term, to compute the total \emph{real emission} NLO corrections. As always, we also want to express the result in such a way that the combination with the virtual IR poles is straightforward. 

Of course, in the sum of the second term of eq. \ref{eq:intFLMm}, there are two disinct classes of corrections: in the case where \(i = a\), the unresolved FS gluon is emmitted by a IS parton. In the second case, it is emitted by another FS gluon. The former is more complicated for two reasons: the initial state parton can be both a gluon or a quark, and the factorization is given by a \emph{rescaled} LO \(F_{LM}\). The integral over the energy \(E_{\mathfrak{m}}\) is not computable because it is process dependent. We must use a strategy similar to what we did for \(pp \to X\) in the previous section.   

We start by considering this last case. The form of the collinear IS limit is universal, 

\begin{align}
    \braket{\bar{S}_{\mathfrak{m}}C_{a\mathfrak{m}}\Delta^{(\mathfrak{m})}F_{LM}(\mathfrak{m})} = \braket{\bar{S}_{\mathfrak{m}}\frac{g^2_{s,b}}{E^2_{\mathfrak{m}}\rho_{\mathfrak{m}a}} \frac{1-z}{z}P_{aa}(z)F_{LM}(z\cdot 1_a, 2_b ; \ldots)}
\end{align}

The angular integral goes as always, 

\begin{align}
    -\frac{[\alpha_s]}{\eps}\bar{S}_{\mathfrak{m}}\left(\frac{2E_1}{\mu}\right)^{-2\eps}\int_{z_{\min}}^1 \frac{dz}{(1-z)^{2\eps}} \frac{F_{LM}(z\cdot1_a,2_b;\ldots)}{z}P_{aa}(z)
\end{align}

We can extract from the kernels \(P_{aa}\) a term that behaves like a delta function in \(z-1\). We write 

\begin{align}
    -\left[ \left( \frac{2E_1}{\mu} \right)^{-2\eps}\frac{\Gamma^2(1-\eps)}{\Gamma(1-2\eps)} \right]P_{aa}(z) = \Gamma_{1a}\delta(1-z) + P_{aa}^{gen}(z)
\end{align}

THe functions \(\Gamma_{ia}\) are called generalized anomalous parton dimension . It is also possible to split \(P_{aa}^{gen}\) as

\begin{align}
    P_{aa}^{gen} = -\left[ \left( \frac{2E_1}{\mu} \right)^{-2\eps}\frac{\Gamma^2(1-\eps)}{\Gamma(1-2\eps)} \right](-\hat{P}^{(0)}_{aa} + \eps P^{fin}_{aa})
\end{align}

where \(\hat{P}^{(0)}\) is the Altarelli-Parisi LO kernel. In this way, the cancellation with the pdfs counterterm (eq. \ref{eq:pdfpartoniccounterterm}) is straightforward. Instead, \(P^{fin}_{aa}\) gives a finite contribution to the cross section as it is a \(o(\eps^0)\) quantity. Finally, the poles proportional to the anomalous dimension of the partons will be canceled with virtual corrections, as they are \(\propto \delta(1-z)\). 

Let's show this explicitly. First, we write the full partonic real cross section contribution as 

\begin{align}
    2sd\hat{\sigma}^R_{ab} = \ldots 
\end{align}

Now we want to combine this with the virtual corrections, that are universal, and can be expressed through Catani's formula: 


\chapter{NNLO QCD corrections}

\section{\(pp \to X\), double real corrections}

The diagrams that contribute to the cross section for \(pp \to X\) at NNLO in QCD can be divided into 3 categories: double real emission, real-virtual, double virtual. Examples can be found in fig \ref{fig:NNLOppX}. Among them, the most complicated, at least for the IR pole structure, are the double real. 

As we did in the NLO case, we subtract first soft and then collinear dirvergences. At the end, we get a fully regolalized and integrable result, plus all the divergent parts, expressed, when possibile, as \(\eps^{-n}\) poles (as before we mork in DR, with \(d = 4-2\eps\)).  

Of course, the poles emerge from integrating the double real corrections over the final gluons phase space, \([dg_4][dg_5]\). This is a two particle phase space, and also the matrix element for the process is more complicated: this explains why the NNLO computation is much harder then the NLO computation. The basic ideas, though, are the same.   

\begin{figure}
    \caption{Examples of diagrams that contribute to \(pp \to X\) at NNLO}
    \label{fig:NNLOppX}
    \begin{multicols}{3}
        \begin{center}
            \begin{tikzpicture}
                \begin{feynman}
                    \vertex (a); 
                    \vertex [right = of a] (b);
                    \vertex [right = of b] (c);  
                    \vertex [below right = of c] (d); 
                    \vertex [right = of d] (e);
                    \vertex [below left = of d] (f);
                    \vertex [left = of f] (g);
                    \vertex [left = of g] (h);
                    \diagram{
                        (a) -- [fermion] (b) -- [fermion] (c) -- [fermion] (d) -- [fermion] (f) -- [fermion] (g) -- [fermion] (h); 
                        (d) -- [boson] (e); 
                        (b) -- [gluon] (g); 
                        (c) -- [gluon] (f); 
                    }; 
                \end{feynman}
            \end{tikzpicture}
        \end{center}
    \columnbreak
    \begin{center}
            \begin{tikzpicture}
                \begin{feynman}
                    \vertex (a); 
                    \vertex [right = of a] (b);
                    \vertex [right = of b] (c);  
                    \vertex [below right = of c] (d); 
                    \vertex [right = of d] (e);
                    \vertex [below left = of d] (f);
                    \vertex [left = of f] (g);
                    \vertex [left = of g] (h);
                    \vertex [right = of c] (l); 
                    \diagram{
                        (a) -- [fermion] (b) -- [fermion] (c) -- [fermion] (d) -- [fermion] (f) -- [fermion] (g) -- [fermion] (h); 
                        (d) -- [boson] (e); 
                        (b) -- [gluon] (g); 
                        (c) -- [gluon] (l); 
                    }; 
                \end{feynman}
            \end{tikzpicture}
        \end{center}
    \columnbreak
        \begin{center}
            \begin{tikzpicture}
                \begin{feynman}
                    \vertex (a); 
                    \vertex [right = of a] (b);
                    \vertex [right = of b] (c);  
                    \vertex [below = of b] (d); 
                    \vertex [below = of d] (e);
                    \vertex [right = of d] (f);
                    \vertex [right = of d] (g);
                    \vertex [left = of e] (h);
                    \vertex [right = of e] (l);
                    \diagram{
                        (a) -- [fermion] (b) -- [fermion] (d) -- [fermion] (e) -- [fermion] (h);
                        (b) -- [gluon] (c); 
                        (d) -- [gluon] (g); 
                        (e) -- [boson] (l); 
                    }; 
                \end{feynman}
            \end{tikzpicture}
        \end{center}
    \end{multicols}
\end{figure}

The full list of limits we have to take to extract the divergences is the following: 

\begin{align}
    S_{ij}, \quad S_i, \quad C_{ij, lm}, \quad C_{ij}
\end{align}

that corrispond to the double-soft, soft, double-colinear, single-colliear limit. In the following sections, we compute all the contributes that emerge when appling the subtraction scheme. 

\subsection{Single-soft}

Let's now focus on the term 

\begin{align}
    \braket{(I-S_{45})S_5 F_{LM}(1, 2, 4, 5)}
\end{align}

To start we have to consider the limit of the amplitued when \(E_5 \to 0\), i.e. \(g_5\) goes soft. We use again Catani's universal formula for the factorization \ref{eq:Catanisoftfact}: 

\begin{align}
    S_5 F_{LM}(1, 2, 4, 5) = -g^2_{s, b} \left[ \sum_{(ij)}S_{ij} (5) \vec{T_i} \cdot \vec{T}_j \right] F_{LM}(1, 2, 4) = -\frac{g_{s,b}^2}{E_5^2}\left[ \sum_{ i < j}\frac{\rho_{ij}}{\rho_{i5}\rho_{j5}} (2\vec{T}_i \cdot \vec{T}_j)\right]F_{LM}(1, 2, 4)
\end{align}

we must evaluate the sum in the squared parenthesis. We remind the reader that in the sum the indexes take values \(i, j = 1, 2, 4, \, i < j\). Color conservation tells us that \(\vec{T}_1 + \vec{T}_2 = \vec{T}_3\). As \(1, 2\) are quarks, \(\vec{T}_1^2 = \vec{T}_2^2 = C_F\), while \(\vec{T}_4^2 = C_A\). We can than manipolate the color conservation equation to read the scalar products: 

\begin{align}
    &(\vec{T}_1 + \vec{T}_2)^2 = \vec{T}_3^2 \implies 2\vec{T}_1\cdot \vec{T}_2 = -(2C_F - C_A) \notag \\ 
    &(\vec{T}_1 - \vec{T}_3)^2 = \vec{T}_2^2 \implies 2\vec{T}_1\cdot \vec{T}_3 = - C_A \notag \\
    &(\vec{T}_2 - \vec{T}_3)^2 = \vec{T}_1^2 \implies 2\vec{T}_2\cdot \vec{T}_3 = - C_A
\end{align}

So the limit reads 

\begin{align}
    S_5 F_{LM}(1, 2, 4, 5) = \frac{g^2_{s, b}}{E_5^2}\left[ (2C_F - C_A)\frac{\rho_{12}}{\rho_{15}\rho_{25}}  + C_A \left( \frac{\rho_{14}}{\rho_{15}\rho_{45}} + \frac{\rho_{24}}{\rho_{25}\rho_{45}} \right)\right] F_{LM}(1,2,4)
\end{align}

All dependence on the momentum of \(g_5\) is in the eikonal factor. We can therefore integrate: the integral is completely analogous to the one to obtain eq. \ref{eq:intsoftlimNLO}. After integrating, we find 

\begin{align}
    &\braket{(I-S_{45})S_5 F_{LM}(1, 2, 4, 5)} = \braket{J_{124} (I-S_{45})F_{LM}(1, 2, 4)}, \notag \\ 
    &\quad J_{124} = \frac{[\alpha_s]}{\eps^2}\left( \frac{E_4}{\mu} \right)^{-2\eps}\left[ (2C_F-C_A)(2\rho_{12})^{-\eps}K_{12} + C_A((2\rho_{14})^{-\eps}K_{14} + (2\rho_{24})^{-\eps}K_{24}) \right]
\end{align}

Unfortunately, this is still not quite the end of the story. We've regulated all divergences that may regard \(g_5\), but we still have to deal with \(g_4\). It can't go soft, because we have removed that possibility. But it can still go collinear to one of the initial state quarks. We therefore have to extract them, using 

\begin{align}
    \braket{J_{124}\bar{S}_4 F_{LM}(1, 2, 4)} = \braket{(\bar{C_{41} + C_{42}})J_{124}\bar{S}_4F_{LM}(1, 2, 4)} + \braket{C_{41} + C_{42}J_{124}\bar{S}_4F_{LM}(1, 2, 4)} 
\end{align}

the first term is free from any divergences, except for the ones that we extracted previously. The second term is what we will be focusing on. We will compute the action of \(C_{41}\) and then deduce the action of \(C_{42}\).  

As we're in the frame of the CM, \(\rho_{12} = 2\), and \(\eta_{12} = \rho_{12}/2 = 1\). It's also straightforward to see that \(C_{14}\rho_{24} = \rho_{14}\). We also remind the reader the following properties of the hypergeometric gauss function, 

\begin{align}
    _2 F_1 (a, b, c, 0) = 1, \quad \lim_{\eta \to 0} \eta^{1+\eps} \, _2 F_1(1, 1, 1-\eps, 1-\eta) = \Gamma(1+\eps)\Gamma(1-\eps) 
\end{align}

Thus the action of \(C_{41}\) on \(J_{124}\) is given by 

\begin{align}
    \lim_{\eta_{14}\to 0} J_{124} = \left( \frac{[\alpha_s]E_4^{-2\eps}}{\eps^2}\frac{\Gamma^2(1-\eps)}{\Gamma(1-2\eps)}\left(2C_F4^{-\eps} + C_A (2\rho_{14})^{-\eps}\Gamma(1-\eps)\Gamma(1+\eps) \right) \right)
\end{align}

Also using the usual universal form of the collinear limit on the amplitude, we get 

\begin{align}
    C_{41}J_{124}\bar{S}_4 F_{LM}(1, 2, 4) = \left( \frac{[\alpha_s]^2E_4^{-2-2\eps}}{4\pi\eps^2\rho_{14}}\frac{\Gamma^2(1-\eps)}{\Gamma(1-2\eps)}\left(2C_F4^{-\eps} + C_A (2\rho_{14})^{-\eps}\Gamma(1-\eps)\Gamma(1+\eps) \right) \right) \times \\ \left( \frac{1-z}{z} \right)P_{qq}(z)F_{LM}(z\cdot 1, 2) 
\end{align}

Now we need to integrate over \([g_4]\). As always the energy integral is process dependent, and cannot be performed. Still, we compute the angular integral and we express the energy integral as an integral in \(z\): the term proportional to \(C_F\) contains the same angular integral that was present in the NLO computation. The term proportional to \(C_A\) is slightly different, and reads 

\begin{align}
    \int \frac{d\cos\theta (1-\cos^2\theta)^{d/2-2}}{\rho_{14} (2\rho_{14})^{-\eps}} = 2^{-2\eps}\int_0^1 d\eta_{14} \eta_{14}^{-1-2\eps}(1-\eta_{14})^{-\eps} = -\frac{2^{-2\eps}}{2\eps}\frac{\Gamma(1-2\eps)\Gamma(1-\eps)}{\Gamma(1-3\eps)}
\end{align}

And the energy integral is therefore given by 

\begin{align}
    -\frac{[\alpha_s]^22^{-2\eps}}{\eps^3}\left[ 2C_F\frac{\Gamma^4(1-\eps)}{\Gamma^2(1-2\eps)} + \frac{C_A}{2}\frac{\Gamma^3(1-\eps)\Gamma(1+\eps)}{\Gamma(1-3\eps)}\right]\int dE_4 E_4^{d-3} E_4^{-2-2\eps} \frac{1-z}{z}P_{qq}(z) F_{LM}(z\cdot 1, 2)
\end{align}

We then switch the integration variable to \(E_4 \to z = 1-E_4/E_1\) (\(\implies E_4 = E_1(1-z)\)), and we obtain 

\begin{align}
    -\frac{[\alpha_s]^2(2E_1/\mu)^{-4\eps}}{\eps^3}\left[ 2C_F\frac{\Gamma^4(1-\eps)}{\Gamma^2(1-2\eps)} + \frac{C_A}{2}\frac{\Gamma^3(1-\eps)\Gamma(1+\eps)}{\Gamma(1-3\eps)}\right]\int_{z_{\min}}^1 dz \frac{(1-z)^{-1-4\eps}}{z}\mathcal{P}_{qq}(z) F_{LM}(z\cdot 1, 2)
\end{align}

Where we abbreviate \(\mathcal{P}_{qq} = \frac{1-z}{z}P_{qq}(z)\). This last equation would be the result of \(\braket{C_{41}J_{124}F_{LM}(1,2, 4)}\). It is not yet what we need, as we also have to insert the operator \(\bar{S}_4\) that makes \(g_4\) not go soft. This is not hard to do, as the soft limit for \(g_4\) is just the limit \(z\to 1\): 

\begin{align}
    \braket{C_{41}J_{124}\bar{S}_4F_{LM}(1,2, 4)} = -\frac{[\alpha_s]^2(2E_1/\mu)^{-4\eps}}{\eps^3}\left[ 2C_F\frac{\Gamma^4(1-\eps)}{\Gamma^2(1-2\eps)} + \frac{C_A}{2}\frac{\Gamma^3(1-\eps)\Gamma(1+\eps)}{\Gamma(1-3\eps)}\right] \times \\ 
    \int_{z_{\min}}^1 dz (1-z)^{-1-4\eps}\hat{\mathcal{P}}^{(-)}_{qq}(z)F_{LM}(z\cdot 1, 2)
\end{align}

the operator \(\hat{\mathcal{P}}^{(-)}_{qq}\) acts on the function of \(z\) \(F_{LM}(z\cdot1, 2)\), as \(\hat{\mathcal{P}}^{(-)}_{qq} f(z) = \mathcal{P}_{qq}f(z) - 2C_F f(1)\). Notice that \(\mathcal{P}_{qq} \xrightarrow{z\to 1} 2C_F\), so \(\mathcal{P}^{(-)}_{qq} f(z)\) goes to zero as \(z \to 1\). 

\subsection{single-colliner (dc partitions)}

We now focus on the term

\begin{align}
    \label{eq:scdc}
    \sum_{(ij) \in \text{dc}}\braket{(I - S_{45})(I-S_5)[C_{4i}[dg_4] + C_{5j}[dg_5]]\omega^{i4, j5}F_{LM}(1, 2, 4, 5)}
\end{align}

starting first from the sector \(14, 25\). Since the action of \(S_{45}\) is the same if there is or there is not \(S_5\), we can drop the first subtraction term. We write then 

\begin{align}
    C_{41}(I-S_5)\omega^{14, 25}F_{LM}(1, 2, 4, 5) = \tilde{\omega}^{14, 25}_{1 \parallel 4}C_{41}(I-S_5)F_{LM}(1, 2, 4, 5)
\end{align}

and using the universal collinear initial state factorization, 

\begin{align}
    C_{41}(I-S_5)\omega^{14, 25}F_{LM}(1, 2, 4, 5) = \tilde{\omega}^{14, 25}_{1 \parallel 4}\frac{g_{s, b}^2}{E_4^2\rho_{14}}\frac{1-z}{z}P_{qq}(z)F_{LM}(z\cdot 1, 2, 5)
\end{align}

where as always, \(z = 1-\frac{E_4}{E_1}\). Before integrating, we notice that in eq. \ref{eq:scdc}, the collinear operator also act on the final state space \([dg_4]\). Before computing the limit, we write the final phase space in a useful way, also absorbing the bare coupling from the collinear limit of the amplitude: 

\begin{align}
    g^2_{s, b} [dg_4]\theta(E_4-E_5) &= g_{s, b}^2 \frac{d^{d-1}P_4}{(2\pi)^{d-1}(2E_4)}\theta(E_4-E_5) = \frac{dE_4 E_4^{d-3}d^{d-2}\Omega}{2 (2\pi)^{d-1}}\theta(E_4-E_5) \notag \\ 
    &= \frac{dE_4 E_4^{d-3} d\eta_{14} 2^{1-2\eps}\eta_{14}^{-\eps}(1-\eta_{14})^{-\eps}g_{s, b}^2 d^{d-3}\Omega}{2(2\pi)^{d-1}}\theta(E_4-E_5) \notag \\ 
    &= [\alpha_s]dE_4 E_4^{d-3} d\eta_{14} 2^{1-2\eps}\eta_{14}^{-\eps}(1-\eta_{14})^{-\eps}\theta(E_4-E_5)
\end{align}

In the last equation we integrated over \(d^{d-3}\Omega\). We also switch \(E_4 \to z\), as we've done many times: 

\begin{align}
    g^2_{s, b} [dg_4]\theta(E_4-E_5) &= 2^{1-2\eps}[\alpha_s]\eta_{14}^{-\eps}(1-\eta_{14})^{-\eps}E_1 dz E_1^{1-2\eps}(1-z)^{1-2\eps}\theta(z(E_5) - z) \notag \\ 
    &= [dg_4]\rho_{14}E_4^2 s^{-\eps} \eta_{14}^{-1-\eps}(1-\eta_{14})^{-\eps} \frac{dz}{(1-z)^{1+2\eps}} \theta(z(E_5) -z)
\end{align}

in the last equation we semplified the expression so it's product with the amplitude is straightforward. Morover, the collinear limit is given by 

\begin{align}
    C_{41}g^2_{s, b} [dg_4]\theta(E_4-E_5) = [dg_4]\rho_{14}E_4^2 s^{-\eps} \eta_{14}^{-1-\eps}\frac{dz}{(1-z)^{1+2\eps}} \theta(z(E_5) -z)
\end{align}

Putting everything together and integrating, we get 

\begin{align}
    \braket{C_{41} [dg_4]\theta(E_4-E_5)F_{LM}(1, 2,4, 5)} &= [\alpha_s] s^{-\eps} \int_0^1 d\eta_{14}\eta_{14}^{-1-\eps}\int_{z_{\min}}^{z_{\max}(E_5)} dz \frac{\mathcal{P}_{qq}(z)\braket{\tilde{\omega}^{14, 25}_{1 \parallel 4}F_{LM}(z\cdot 1, 2, 5)}}{(1-z)^{1+2\eps}} \notag \\ 
    &= -\frac{[\alpha_s] s^{-\eps}}{\eps}\int_{z_{\min}}^{z_{\max}(E_5)} dz \mathcal{P}_{qq}(z)\frac{\braket{\tilde{\omega}^{14, 25}_{1 \parallel 4}F_{LM}(z\cdot 1, 2, 5)}}{(1-z)^{1+2\eps}}
\end{align}

Now, we need to consider the term with \(C_{52}\). Following the same steps as before, we get 

\begin{align}
    \braket{(I-S{45})(I-S_5)C_{52}w^{14, 25}[dg_5]F_{LM}(1, 2, 4, 5)} = -\frac{[\alpha_s]s^{-\eps}}{\eps}\int_{z_{\min}(E_4)}^1 \frac{dz\hat{\mathcal{P}}^{(-)}_{qq}}{(1-z)^{1+2\eps}} \braket{\tilde{w}^{14, 25}_{5\parallel 2}F_{LM}(1, z\cdot 2, 4)}
\end{align}

here of course \(z = 1-E_5/E_2\). The result slightly differs becouse 

\begin{enumerate}
    \item \(E_5\) must be smaller than \(E_4\), so the integration over \(z\) has to be performed up to \(z=1\). 
    \item The operator \(I-S_5\) now acts on the gluon that is collinear: we therefore treat it introducing \(\hat{\mathcal{P}}^{(-)}_{qq}\), exactly like in the single-soft term. 
\end{enumerate}

Summing the two results, we get the sector \(14, 25\). The other double-collinear sector is the \(15, 24\), and its result is the same, with the obvious replacements. To conclude this computation we must add up this four contributions. 

\subsection{single collinear - (tc partitions)}

We now want to compute 

\begin{align}
    \sum_{i \in \text{tc}} \braket{ (I-S_{45}(I-S_5))[\theta^{(a)}C_{5i} + \theta^{(b)}C_{45} + \theta^{(c)}C_{4i} + \theta^{(d)}C_{45}] [dg_4][dg_5] w^{i4, i5}F_{LM}(1, 2, 4, 5)}
\end{align}

we start by

\begin{align}
    \braket{(I-S_{45})(I-S_5)\theta^{(a)}C_{51}[dg_5]w^{14, 15}F_{LM}(1, 2, 4, 5)}
\end{align}

This is similar to the previous computation. The theta factor \(\theta^{(a)} = \theta(\eta_{14}/2-\eta_{15})\) modifies the extrems of integration in the angular integral of \(g_5\), in the following way: 

\begin{align}
    \int_0^1 d\eta_{15} \frac{\theta^{(a)}}{\eta_{15}^{1+\eps}} = \int_0^{\eta_{14}/2} \frac{d\eta_{15}}{\eta_{15}^{1+\eps}} = -\frac{(\eta_{14}/2)^{-\eps}}{\eps}
\end{align}

the rest is analogous. Since \(g_5\) is both collinear and hard, we get

\begin{align}
    \braket{(I-S_{45})(I-S_5)\theta^{(a)}C_{51}[dg_5]w^{14, 15}F_{LM}(1, 2, 4, 5)} = -\frac{[\alpha_s]s^{-\eps}}{\eps} \int_{z_{\min(E_4)}}^1 dz \frac{\hat{\mathcal{P}}_{qq}^{(-)}}{(1-z)^{1+2\eps}}\braket{\tilde{w}^{14, 15}_{1 \parallel 5} (\rho_{14}/4)^{-\eps}F_{LM}(z\cdot 1, 2, 4)}
\end{align}

We treat the term 

\begin{align}
    \braket{(I-S_{45})(I-S_5)\theta^{(c)}C_{41}[dg_4]w^{14, 15}F_{LM}(1, 2, 4, 5)}
\end{align}

because it's very similar. This time it's \(g_4\) that goes collinear, so the result is given by 

\begin{align}
    \braket{(I-S_{45})(I-S_5)\theta^{(c)}C_{41}[dg_4]w^{14, 15}F_{LM}(1, 2, 4, 5)}  = -\frac{[\alpha_s]s^{-\eps}}{\eps} \int_{z_{\min(E_4)}}^{z_{\max}(E_5)} dz \frac{\hat{\mathcal{P}}_{qq}^{(-)}}{(1-z)^{1+2\eps}}\braket{\tilde{w}^{14, 15}_{1 \parallel 4} (\rho_{15}/4)^{-\eps}F_{LM}(z\cdot 1, 2, 5)}
\end{align}

Instead, the sectors \((b), (d)\) are very different, as they are the sectors in wich \(g_4, g_5\) become collinear between each other. We start with 

\begin{align}
    \braket{(I-S_{45})(I-S_5)\theta^{b, d}C_{45}w^{14, 15}[dg_5]F_{LM}(1, 2, 4, 5)}
\end{align}

the collinear limit on the amplitude is the standard final state collinear limit: 

\begin{align}
    C_{45}F_{LM}(1, 2, 4, 5) &= \frac{g^2_{s, b}}{E_5^2 \rho_{45}}\frac{E_5}{E_4}P_{gg, \mu\nu}(z)F^{\mu\nu}_{LM}(1, 2, 4+5) \notag \\ 
    &= \frac{g^2_{s, b}}{E_5^2 \rho_{45}}\frac{z}{1-z}\left[P^{(0)}_{gg}(z)F_{LM}(1, 2, 45) + P_{gg}^\perp(z)k_{\perp, \mu}k_{\perp, \nu}F^{\mu\nu}_{LM}(1, 2, 45) \right]
\end{align}

with 

\begin{align}
    P^{(0)}_{gg} = 2C_A \left( \frac{z}{1-z} + \frac{1-z}{z} \right), \quad P^\perp_{gg}(z) = 4C_A(1-\eps)z(1-z)
\end{align}

here \(z\) is the fraction of the total momentum, \(p_{45} = p_4 + p_5\) that is carried by \(p_5\). Instead, \(k_\perp\) is a trasverse vector. It comes from the following decomposition 

\begin{align}
    p_5 = \alpha p_4 + \beta \bar{p}_4 + k_\perp 
\end{align}

where \(\bar{p}_4 = (p_4^0, -\vec{p}_4)\), and we call it trasverse because \(k_\perp \cdot p_4 = k_\perp \cdot \bar{p}_4 = 0\). 

In order to proceed, we need an explicit parametrization of these vectors and the final phase space. Consider then 

\begin{align}
    [dg_4][dg_5]\theta(E_{\max}-E_4)\theta(E_4 - E_5)
\end{align}

We can take \(E_4 = x_1 E_{\max}\), \(E_5 = x_2 E_4 = x_2 x_1 E_{\max}\), with \(x_1, x_2 \in [0, 1]\). In this way, we write 

\begin{align}
    [dg_4][dg_5] &= dE_4 E_4^{d-3}dE_5 E_5^{d-3} d\Omega_{45}, \quad d\Omega_{45} = \frac{d^{d-2}\Omega_4}{2(2\pi)^{d-1}}\frac{d^{d-2}\Omega_5}{2(2\pi)^{d-1}} \notag \\ 
    &= E_{\max}^{4-4\eps}(dx_1 x_1^{3-2\eps})(dx_2 x_2^{1-2\eps})
\end{align}















\end{document}

 
